\subsection{Equivariant functions}

In this subsection, we fix a closed subgroup H of G. We denote by \(\varpi\) the canonical projection of G onto G/H.

In addition to the Haar measure \(\mu\) on G, we fix a left Haar measure \(\beta\) on H.

By INT, VII, p. 56, § 2, no. 5, th. 2, there exists a continuous function \(\kappa\colon\mathbf{G}\to\mathbf{R}_{+}^{*}\) such that \(\kappa(xh)=\Delta_{\mathrm{H}}(h)\Delta_{\mathrm{G}}(h)^{-1}\kappa(x)\) for every \((x,h)\in\mathbf{G}\times\mathbf{H}\) . We fix such a function \(\kappa\) . We denote \(\nu\) the measure \((\kappa\cdot\mu)/\beta\) on G/H; by loc. cit., it is a non-zero positive measure quasi-invariant under G. Its support is equal to G/H (INT,

VII, p. 10, § 1, no. 1). By INT, VII, p. 43, § 2, no. 2, prop. 4, the measure \(\nu\) is the unique measure on G/H such that the measure \(\nu^{\sharp}\) on G is equal to \(\kappa \cdot \mu\) . We equip the space G/H with the measure \(\nu\) (so that we write, for example, \(\mathcal{L}^{p}(\mathrm{G / H}) = \mathcal{L}^{p}(\mathrm{G / H},\nu)\) ).

 We say that a set \(S \subset G\) is negligible modulo H if \(\varpi(S)\) is \(\nu\) -negligible. This condition does not depend on the choice of Haar measures on G and on H. It implies that S is locally \(\mu\) -negligible (INT, VII, p. 47, § 2, no. 3, prop. 6, a)). We say that a property of an element of G is true almost everywhere modulo H if the set of elements of G for which this property does not hold is negligible modulo H.

Let \(\pi\) a unitary representation of H in a Hilbert space E. We denote by \(\mathcal{F}_{\pi}(\mathrm{G})\) the vector space of functions \(f\) on G with values in E such that \(f(xh) = \pi(h)f(x)\) for all \((x,h) \in \mathrm{G} \times \mathrm{H}\) . We say that the elements of \(\mathcal{F}_{\pi}(\mathrm{G})\) are the \(\pi\) -equivariant functions on G.

The space \(\mathcal{F}_{1}(\mathrm{G})\) associated with the trivial representation of G on C is identified with the space \(\mathcal{F}(\mathrm{G / H})\) of complex-valued functions on G/H by the map \(f\mapsto f\circ \varpi\) of \(\mathcal{F}(\mathrm{G / H})\) in \(\mathcal{F}_1(\mathrm{G})\) .

For any function \(f\) in \(\mathcal{F}_{\pi}(\mathrm{G})\) , the function \(\|f\|\) belongs to \(\mathcal{F}_{1}(\mathrm{G})\) since \(\pi\) is unitary. This allows us to identify \(\|f\|\) with a function on \(\mathrm{G/H}\) , and we shall write, for example, \(\|f(x\mathrm{H})\|\) for the value of this function at an element \(x\mathrm{H}\) of \(\mathrm{G/H}\) .

A function \(f\) of \(\mathcal{F}_{\pi}(\mathrm{G})\) will be said to vanish outside a compact set modulo H if \(\|f\|\) vanishes outside a compact subset of G/H. This is equivalent to saying that there exists a compact subset K of G such that \(f\) vanishes outside \(\mathrm{K} \cdot \mathrm{H}\) (TG, III, p. 33, prop. 10).

We denote \(\mathcal{K}_{\pi}(\mathrm{G})\) the space of continuous functions on G belonging to \(\mathcal{F}_{\pi}(\mathrm{G})\) which have compact support modulo H.

A space analogous to \(\mathcal{H}_{\pi}(\mathrm{G})\) was defined in INT, VII, p. 39, § 2, no. 1, when \(\pi\) is a continuous homomorphism of H into \(\mathbf{R}_{+}^{*}\) .

Let \(p \in [1, +\infty[\) . For every \(f \in \mathcal{F}_{\pi}(\mathrm{G})\) , we denote

\[
\mathrm{N} _ {p} (f) = \left(\int_ {\mathrm{G/H}} ^ {*} \| f \| ^ {p} d \nu\right) ^ {1 / p}.
\]

It is a real number or \(+\infty\) . We denote \(\mathcal{F}_{\pi}^{p}(\mathrm{G},\nu)\) , or simply \(\mathcal{F}_{\pi}^{p}(\mathrm{G})\) , the subspace of functions \(f\in \mathcal{F}_{\pi}(\mathrm{G})\) such that \(\mathrm{N}_p(f)\)

is finite. The space \(\mathcal{F}_{\pi}^{p}(G)\) equipped with the map \(N_{p}\) is a semi-normed space.

The space \(\mathcal{K}_{\pi}(\mathrm{G})\) is contained in \(\mathcal{F}_{\pi}^{p}(\mathrm{G})\) . Its closure in \(\mathcal{F}_{\pi}^{p}(\mathrm{G})\) is denoted by \(\mathcal{L}_{\pi}^{p}(\mathrm{G},\nu)\) , or simply \(\mathcal{L}_{\pi}^{p}(\mathrm{G})\) ; we say that it is the space of \(\pi\) -equivariant functions on G of power \(p^{e}\) integrable modulo H.

The assertions of the following proposition, when \(\pi\) is the trivial representation of dimension 1, are consequences of INT, IV, p. 128, § 3, no. 3, prop. 6 and p. 131, § 3, no. 4, th. 3.

PROPOSITION 3. --- a) Let \((f_n)_{n \in \mathbf{N}}\) a sequence in \(\mathcal{F}_\pi^p(\mathrm{G})\) such that the series

\[
\sum_ {n = 0} ^ {+ \infty} \mathrm{N} _ {p} (f _ {n})
\]

converges. Then the series with general term \(f_n(g)\) is absolutely convergent for \(g\) outside a set T which is negligible modulo H. Let \(f\) be the function on G with values in E which is equal to the sum of this series on G---T and equal to 0 on T. We then have \(f \in \mathcal{F}_\pi^p(\mathrm{G})\) and the series with general term \(f_n\) converges to \(f\) in the space \(\mathcal{F}_\pi^p(\mathrm{G})\) ;

\begin{enumerate}
\def\labelenumi{\alph{enumi})}
\setcounter{enumi}{1}
\item
  Let \((f_n)\) a sequence converging in \(\mathcal{L}_{\pi}^{p}(\mathrm{G})\) to a function \(f\) . There exists a sequence \((f_{n_k})\) extracted from \((f_n)\) such that \(f_{n_k}(g)\) converges to \(f(g)\) for all \(g\) outside a set negligible modulo H;
\item
  The semi-normed spaces \(\mathcal{F}_{\pi}^{p}(\mathrm{G})\) and \(\mathcal{L}_{\pi}^{p}(\mathrm{G})\) are complete.
\end{enumerate}

In the assertion \(a\) ), saying that the series with general term \(f_n\) converges to \(f\) in the space \(\mathcal{F}_{\pi}^{p}(\mathrm{G})\) means that the sequence of partial sums \(f_0 + \cdots + f_n\) converges to \(f\) in \(\mathcal{F}_{\pi}^{p}(\mathrm{G})\) . One then also says that \(f\) is a sum of this series.

Let us prove \(a\) ). By prop. 6 of INT, IV, p. 128, § 3, no. 3, there exists a set \(\nu\) -negligible S ⊂ G/H such that the series with general term \(\|f_n(gH)\|\) converges absolutely for gH \(\notin\) . S. Moreover, the function h which is equal to the sum of this series for gH \(\notin\) . S and which is zero for gH \(\in\) . S satisfies N \(_p(h)\) \textless{} + \(\infty\) .

The set T = \(\overline{\varpi}(S)\) is negligible modulo H. For every \(g \notin T\) , the series with general term \(f_{n}(g)\) is absolutely convergent in E. Let us define \(f(g) = \sum f_{n}(g)\) for \(g \notin T\) and \(f(g) = 0\) otherwise. We have \(f \in \mathcal{F}_{\pi}(G)\) . Note that \(\|f(gH)\| \leqslant h(gH)\) for all \(g \in G\) , whence \(N_{p}(f) \leqslant\)

\(N_{p}(h) < +\infty\) . We therefore have \(f \in \mathcal{F}_{\pi}^{p}(G)\) . Similarly, it follows that

\[
\mathrm{N} _ {p} \Big (f - \sum_ {n = 0} ^ {k} f _ {n} \Big) \leqslant \sum_ {n = k + 1} ^ {+ \infty} \mathrm{N} _ {p} (f _ {n})
\]

for all \(k \in \mathbf{N}\) , hence the series \(\sum f_n\) converges to \(f\) in \(\mathcal{F}_{\pi}^{p}(\mathrm{G})\) . The assertion \(a)\) is proved.

Let us prove \(b\) ). The sequence \((\|f_n - f\|)_{n \in \mathbf{N}}\) converges to 0 in the space \(\mathcal{L}^p(\mathrm{G}/\mathrm{H})\) . By INT, IV, p. 131, § 4, no. 3, th. 3, there exists a subsequence \((\|f_{n_k} - f\|)_{k \in \mathbf{N}}\) which converges \(\nu\) -almost everywhere to 0. The sequence \((f_{n_k}(g))_{k \in \mathbf{N}}\) then converges to \(f(g)\) for all \(g\) outside a set negligible modulo H.

Finally let us prove \(c\) ). Let \((f_n)_{n \in \mathbf{N}}\) a Cauchy sequence in \(\mathcal{F}_\pi^p(\mathrm{G})\) . There exists a strictly increasing sequence of integers \((n_k)_{k \in \mathbf{N}}\) such that \(\mathrm{N}_p(f_{n_{k+1}} - f_{n_k}) \leqslant 2^{-k}\) for all \(k \in \mathbf{N}\) . For every \(k \in \mathbf{N}\) set \(h_k = f_{n_{k+1}} - f_{n_k} \in \mathcal{F}_\pi^p(\mathrm{G})\) . By \(a\) ), the series with general term \(h_k\) converges in \(\mathcal{F}_\pi^p(\mathrm{G})\) ; denote it \(h\) its sum. For every \(\ell \in \mathbf{N}\) , it follows that

\[
f _ {n _ {0}} + \sum_ {k = 0} ^ {\ell} h _ {k} = f _ {n _ {\ell + 1}},
\]

hence \(f_{n_{0}} + h\) is a cluster point of the sequence \((f_{n})\) . The latter is therefore convergent (TG, II, p. 14, cor. 2 of prop. 5). Thus the space \(\mathcal{F}_{\pi}^{p}(\mathrm{G})\) is complete; since the space \(\mathcal{L}_{\pi}^{p}(\mathrm{G})\) is closed in \(\mathcal{F}_{\pi}^{p}(\mathrm{G})\) , it is also complete.

COROLLARY. --- Let \((f_{n})_{n\in\mathbf{N}}\) a Cauchy sequence in \(\mathcal{L}_{\pi}^{p}(\mathrm{G})\) and let \(f\in\mathcal{F}_{\pi}(\mathrm{G})\) such that \(f_{n}(g)\) converges to \(f(g)\) almost everywhere modulo H. Then \(f\in\mathcal{L}_{\pi}^{p}(\mathrm{G})\) and \((f_{n})_{n\in\mathbf{N}}\) converges to f in \(\mathcal{L}_{\pi}^{p}(\mathrm{G})\) .

Indeed, the function \(f\) coincides almost everywhere modulo H with the limit of the sequence \((f_n)\) in \(\mathcal{L}_{\pi}^{p}(\mathrm{G})\) .

We denote \(\mathrm{L}_{\pi}^{p}(\mathrm{G})\) the separated normed topological vector space associated with the seminormed space \(\mathcal{L}_{\pi}^{p}(\mathrm{G})\) ; it is a Banach space.

Let \(f\) a function on \(\mathrm{G}\) with values in \(\mathrm{E}\) . We denote \(\mathrm{S}_{f}\) the set of \(x\in \mathrm{G}\) such that the function \(h\mapsto f(xh)\) onto \(\mathrm{H}\) does not belong to \(\mathcal{L}_{\mathrm{E}}^{1}(\mathrm{H})\) . We have \(\mathrm{S}_f\cdot h = \mathrm{S}_f\) for all \(h\in \mathrm{H}\) .

Let \(x \in \mathrm{G} - \mathrm{S}_f\) . Since \(\pi\) is a unitary representation, the function \(h \mapsto \pi(h)^* f(xh)\) is measurable (as a composition of the map \(h \mapsto (h, f(xh))\) from H to \(\mathrm{H} \times \mathrm{E}\) which is then measurable and of the continuous map \((g,x)\mapsto\pi(g)^{*}x\) from H × E to E, cf. INT, IV, p. 174, § 5, no. 3, th. 1), and it is integrable on H.

One defines a function \(f^{\pi}\) on G by setting

\[
f ^ {\pi} (x) = \int_ {\mathrm{H}} \pi (h) ^ {*} f (x h) d \beta (h)\tag{3}
\]

for \(x \in \mathrm{G} - \mathrm{S}_f\) and \(f^\pi(x) = 0\) if \(x \in \mathrm{S}_f\) .

PROPOSITION 4. --- Let \(f \in \mathcal{L}_{\mathrm{E}}^{1}(\mathrm{G})\) .

\begin{enumerate}
\def\labelenumi{\alph{enumi})}
\item
  The set \(\mathrm{S}_f\) is negligible modulo H and \(f^\pi \in \mathcal{F}_\pi(\mathrm{G})\) ;
\item
  Let C be a compact subset of G. If f is continuous and supported in C, then \(S_{f}\) is empty, the function \(f^{\pi}\) belongs to \(\mathcal{H}_{\pi}(\mathrm{G})\) and its support is contained in \(C \cdot H\) .
\end{enumerate}

The first part of \(a\) ) follows from INT, VII, p. 57, § 2, no. 5, c). Let \(w \in H\) . If \(x \in G - S_f\) , then \(xw \in G - S_f\) and

\[
\begin{array}{c} f ^ {\pi} (x w) = \int_ {\mathrm{H}} \pi (h) ^ {*} f (x w h) d \beta (h) \\ = \int_ {\mathrm{H}} \pi (w ^ {- 1} y) ^ {*} f (x y) d \beta (y) = \pi (w) f ^ {\pi} (x), \end{array}
\]

whereas if \(x \in S_{f}\) , then \(f^{\pi}(xw) = 0 = \pi(w)^{*}f^{\pi}(x)\) . The function \(f^{\pi}\) therefore belongs to \(\mathcal{F}_{\pi}(\mathrm{G})\) .

Suppose now that \(f\) is continuous and supported in C. For every \(x \in G\) , the map \(h \mapsto f(xh)\) then belongs to \(\mathcal{K}(H)\) , hence is integrable, which proves that \(S_f = \varnothing\) .

Let us prove that \(f^{\pi}\) is continuous. Let \(x \in G\) . Let U be a relatively compact open neighborhood of \(x\) in G.

For every \(y \in \mathrm{U}\) , we have

\[
\begin{array}{l} \| f ^ {\pi} (y) - f ^ {\pi} (x) \| \leqslant \int_ {\mathrm{H}} \| f (y h) - f (x h) \|   d \beta (h) \\ \qquad = \int_ {\mathrm{H} \cap (y ^ {- 1} \mathrm{C} \cup x ^ {- 1} \mathrm{C})} \| f (y h) - f (x h) \|   d \beta (h) \\ \qquad \leqslant \beta (\mathrm{H} \cap \mathrm{U} ^ {- 1} \mathrm{C}) \sup _ {h \in \mathrm{U} ^ {- 1} \mathrm{C}} \| f (y h) - f (x h) \|, \end{array}
\]

and the continuity of \(f^{\pi}\) then follows from the uniform continuity of f on G.

Finally, if \(x \in G\) satisfies \(f^{\pi}(x) \neq 0\) , there exist \(h \in H\) such that \(f(xh) \neq 0\) thus xh belongs to C and x belongs to C·H. Since C·H is closed in G, we conclude that the support of \(f^{\pi}\) is contained in C·H.

 Let C be a compact subset of G. Let \(u \in \mathcal{K}_{+}(G)\) a function such that \(u(x) > 0\) for all \(x \in C\) . Let v be the function on G defined by

\[
v (x) = \int_ {\mathrm{H}} u (x h) d \beta (h)\tag{4}
\]

for all \(x \in G\) ; with the previous notation, we have \(v = u^{1}\) , corresponding to the trivial 1-dimensional representation of H. The function v is continuous and positive; it belongs to \(\mathcal{F}_{1}(G)\) , its support is contained in \(\operatorname{Supp}(u) \cdot H\) and we have

\[
\inf _ {x \in \mathrm{C} \cdot \mathrm{H}} v (x) > 0\tag{5}
\]

(INT, VII, p. 39--40, § 2, no. 1, prop. 1 and lemma 1, a)).

Lemma 5. --- Let \(f \in \mathcal{F}_{\pi}(\mathrm{G})\) a function \(\mu\) -measurable and zero outside of \(\mathrm{C} \cdot \mathrm{H}\) such that the function \(\|f\|\) is \(\nu\) -integrable on \(\mathrm{G}/\mathrm{H}\) . Let s be the function on \(\mathrm{G}\) with values in E such that

\[
 s (x) = \left\{ \begin{array}{l l} v (x) ^ {- 1} f (x) & \text{if } x \in \mathrm{C} \cdot \mathrm{H} \\ 0 & \text{if } x \in \mathrm{G} - \mathrm{C} \cdot \mathrm{H}. \end{array} \right.
\]

\begin{enumerate}
\def\labelenumi{\alph{enumi})}
\item
  The function s is \(\mu\) -measurable; it belongs to \(\mathcal{F}_{\pi}(\mathrm{G})\) and is zero outside of C·H;
\item
  The function us belongs to \(\mathcal{L}^1 (\mathrm{G})\) and \((us)^{\pi} = f\) almost everywhere modulo H.
\end{enumerate}

The function \(s\) vanishes outside C \(\cdot\) . H by definition; it is \(\mu\) -measurable since the function \(f\) is \(\mu\) -measurable and \(v(x) > 0\) for all \(x \in \mathrm{C} \cdot \mathrm{H}\) (INT, IV, p. 193, § 5, no. 10, prop. 16). We have \(s \in \mathcal{F}_{\pi}(\mathrm{G})\) because \(v \in \mathcal{F}_1(\mathrm{G})\) .

The function \(f\) is locally \(\mu\) -integrable, since \(\|f\|\) is a function \(\nu\) -integrable (INT, VII, p. 47, § 2, no. 3, prop. 6, c), noting that the measure \(\nu^{\sharp} = \kappa \cdot \mu\) is equivalent to \(\mu\) ), hence the function \(s\) is likewise by formula (5). The function \(us\) is measurable and zero outside a compact subset of G; since \(u\) is bounded, it follows that \(us\) is \(\mu\) -integrable.

For every \(x\in \mathrm{G - S}_{us}\) , we have

\[
\begin{array}{c} (u s) ^ {\pi} (x) = \int_ {\mathrm{H}} \pi (h) ^ {*} (u (x h) s (x h)) d \beta (h) \\ = \int_ {\mathrm{H}} u (x h) s (x) d \beta (h) = v (x) s (x), \end{array}
\]

since \(s(xh) = \pi(h)s(x)\) ; the last assertion follows from prop. 4, a).

When \(H = \{e\}\) and \(\pi\) is of dimension 1, the following proposition is none other than INT, IV, p. 184, § 5, no. 6, th. 5.

PROPOSITION 5. --- Let \(p \in [1, +\infty[\) . The space \(\mathcal{L}_{\pi}^{p}(\mathrm{G})\) is the space of functions \(f \in \mathcal{F}_{\pi}(\mathrm{G})\) such that \(f\) is \(\mu\) -measurable and that the function \(\|f\|\) belongs to \(\mathcal{L}^{p}(\mathrm{G}/\mathrm{H})\) .

Let \(f \in \mathcal{L}_{\pi}^{p}(G)\) . It is the limit in \(\mathcal{L}_{\pi}^{p}(G)\) of a sequence of elements of \(\mathcal{K}_{\pi}(G)\) . By prop. 3, b) and Egoroff's theorem (INT, IV, p. 175, § 5, no. 4, th. 2), the function \(f\) is therefore \(\mu\) -measurable; consequently, the function \(\|f\|\) on G/H is \(\nu\) -measurable (INT, VII, p. 47, § 2, no. 3, prop. 6, b)). Since \(N_{p}(f)\) is finite, the function \(\|f\|\) belongs to \(\mathcal{L}^{p}(G/H)\) (INT, IV, p. 184, § 5, no. 6, th. 5).

 Let us prove the converse assertion. Let \(f\) a function \(\mu\) -measurable belonging to \(\mathcal{F}_{\pi}(\mathrm{G})\) such that \(\|f\|\in\mathcal{L}^{p}(\mathrm{G}/\mathrm{H})\) . Let \(\varepsilon>0\) . Let us prove that there exists \(\widetilde{f}\in\mathcal{K}_{\pi}(\mathrm{G})\) such that

\[
\mathrm{N} _ {p} (f - \widetilde {f}) ^ {p} = \int_ {\mathrm{G/H}} ^ {*} \| f - \widetilde {f} \| ^ {p} d \nu \leqslant \varepsilon ,
\]

which will conclude the proof.

 Suppose first that there exists a compact subset C of G such that \(f\) vanishes outside \(\mathrm{C} \cdot \mathrm{H}\) . Let \(u \in \mathcal{K}_{+}(\mathrm{G})\) a function such that \(u(x) > 0\) for all \(x \in \mathrm{C}\) and let us define \(v = u^{1}\) as above. Denote \(\varphi\) the characteristic function of the support of \(u\) .

Let q be the conjugate exponent of p. If p = 1, let w be the constant function on G equal to \(\sup_{x \in G} u(x)\) . If p \textgreater{} 1, we define

\[
w (x) = \left(\int_ {\mathrm{H}} u (x h) ^ {q} d \beta (h)\right) ^ {1 / q}
\]

 for all \(x \in G\) . In all cases, the function w is continuous and positive; it belongs to \(\mathcal{K}_{1}(G)\) (Prop. 4 applied to the trivial representation of dimension 1 and to the function \(u^{q}\) ) hence it is bounded on G. We set \(W = \sup_{x \in G} w(x)\) .

 Let \(s\) the function on \(\mathbf{G}\) with values in \(\mathbf{E}\) defined by Lemma 5 applied to \(f\) . Set \(g = \varphi s\) . The function \(g\) is \(\mu\) -measurable, and satisfies \(\| g \| \leqslant \| f \| / \inf_{x \in \mathrm{C} \cdot \mathrm{H}} v(x)\) . Since \(g\) is zero outside the support of \(u\) and that \(\kappa\) is continuous, we have \(g \in \mathcal{L}_{\mathrm{E}}^{p}(\mathrm{G}, \kappa \cdot \mu)\) . Let \(\widetilde{g}\) a function

in \(\mathcal{K}_{\mathrm{E}}(\mathrm{G})\) such that

\[
\int_ {\mathrm{G}} ^ {*} \| g - \widetilde {g} \| ^ {p} \kappa d \mu \leqslant \frac {\varepsilon}{\mathrm{W} ^ {p}}.
\]

We have us = ug, hence \(f = (us)^{\pi} = (ug)^{\pi}\) almost everywhere modulo H (lemma 5, b)). Let \(\widetilde{f} = (u\widetilde{g})^{\pi}\) ; we have \(\widetilde{f} \in \mathcal{K}_{\pi}(\mathrm{G})\) (prop. 4, b)). For every \(x \in G - S_{ug}\) , we obtain

\[
\| (u g) ^ {\pi} (x) - (u \widetilde {g}) ^ {\pi} (x) \| \leqslant \int_ {\mathrm{H}} ^ {*} u (x h) \| g (x h) - \widetilde {g} (x h) \| d \beta (h)
\]

whence

\[
\| (u g) ^ {\pi} (x) - (u \widetilde {g}) ^ {\pi} (x) \| ^ {p} \leqslant w (x) ^ {p} \int_ {\mathrm{H}} ^ {*} \| g (x h) - \widetilde {g} (x h) \| ^ {p} d \beta (h),
\]

by Hölder's inequality in the case \( p \) \textgreater{} 1. Since \(S_{ug}\) is negligible modulo H, it then follows that

\[
\begin{array}{c} \int_ {\mathrm{G/H}} ^ {*} \| f - \widetilde {f} \| ^ {p} d \nu \leqslant \mathrm{W} ^ {p} \int_ {\mathrm{G/H}} ^ {*} \Big (\int_ {\mathrm{H}} ^ {*} \| g (x h) - \widetilde {g} (x h) \| d \beta (h) \Big) ^ {p} d \nu (x \mathrm{H}) \\ = \mathrm{W} ^ {p} \int_ {\mathrm{G}} ^ {*} \| g - \widetilde {g} \| ^ {p} \kappa d \mu \leqslant \varepsilon \end{array}
\]

according to INT, VII, p. 46, § 2, no. 3, prop. 5, b), which is applicable because \(g - \widetilde{g}\) vanishes outside a compact subset of G. This implies the required property when \(f\) is zero outside a compact subset modulo H.

Now consider the general case. Since \(\|f\|\in\mathcal{L}^{p}(G/H)\) by hypothesis, there exists a compact subset L of G/H such that

\[
\int_ {(\mathrm{G} / \mathrm{H}) - \mathrm{L}} ^ {*} \| f \| ^ {p} d \nu \leqslant \frac {\varepsilon}{2}
\]

(cf. INT, IV, p. 152, § 4, no. 6, th. 4). Let \(\varphi_{\mathrm{L}}\) the characteristic function of L and set \(f_{\mathrm{L}} = (\varphi_{\mathrm{L}} \circ \varpi)f\) . We have

\[
\int_ {\mathrm{G} / \mathrm{H}} ^ {*} \| f - f _ {\mathrm{L}} \| ^ {p} d \nu = \int_ {(\mathrm{G} / \mathrm{H}) - \mathrm{L}} ^ {*} \| f \| ^ {p} d \nu \leqslant \frac {\varepsilon}{2}.
\]

The function \(f_{\mathrm{L}}\) is \(\mu\) -measurable and zero outside a compact set modulo H. It belongs to \(\mathcal{L}_{\pi}^{p}(\mathrm{G})\) therefore, by the preceding case, there exists \(\widetilde{f} \in \mathcal{K}_{\pi}(\mathrm{G})\) such that \(\mathrm{N}_p(f_{\mathrm{L}} - \widetilde{f}) \leqslant (\frac{\varepsilon}{2})^{1/p}\) , whence

\[
\int_ {\mathrm{G/H}} ^ {*} \| f - \widetilde {f} \| ^ {p} d \nu \leqslant \varepsilon ,
\]

as desired.

 Let us consider the case \(p = 2\) . For \(f_1\) and \(f_2\) in \(\mathcal{F}_{\pi}(\mathrm{G})\) , the function \(x \mapsto \langle f_1(x) | f_2(x) \rangle\) belongs to \(\mathcal{F}_1(\mathrm{G})\) since the representation \(\pi\) is unitary, and therefore defines by passing to the quotient a function on \(\mathrm{G / H}\) which we identify as before with \(\langle f_1 | f_2 \rangle\) . We have the bound \(|\langle f_1 | f_2 \rangle| \leqslant \| f_1 \| \| f_2 \|\) in \(\mathcal{F}(\mathrm{G / H})\) .

If \(f_1\) and \(f_2\) belong to \(\mathcal{L}_\pi^2(\mathrm{G})\) , then the function \(\langle f_1 | f_2 \rangle\) belongs to \(\mathcal{L}_1^1(\mathrm{G})\) . In particular, it is integrable on \(\mathrm{G/H}\) . The map

\[
(f _ {1}, f _ {2}) \mapsto \int_ {\mathrm{G} / \mathrm{H}} \langle f _ {1} | f _ {2} \rangle d \nu
\]

is a positive hermitian form on \(\mathcal{L}_{\pi}^{2}(\mathrm{G})\) ; the associated semi-norm is the semi-norm \(N_{2}\) . In particular, the space \(\mathcal{L}_{\pi}^{2}(\mathrm{G})\) is a pre-Hilbert space and \(\mathrm{L}_{\pi}^{2}(\mathrm{G})\) is the Hilbert space associated with \(\mathcal{L}_{\pi}^{2}(\mathrm{G})\) .

\subsection{Induced representations}

We keep the notations and conventions of the preceding number concerning the measures \(\beta\) on H and \(\nu\) on G/H, as well as the function \(\kappa\colon\mathrm{G}\to\mathbf{R}_{+}^{*}\) .

There exists a continuous function \(\eta\) of \(\mathrm{G} \times \mathrm{G} / \mathrm{H}\) in \(\mathbf{R}_{+}^{*}\) such that

\[
\eta (x, y \mathrm{H}) = \frac {\kappa (x y)}{\kappa (x)}
\]

for all \((x,y)\in\mathrm{G}\times\mathrm{G}\) , and \(\boldsymbol{\gamma}_{\mathrm{G}/\mathrm{H}}(x)\nu=(y\mapsto\eta(x^{-1},y))\cdot\nu\) for \(x\in G\) (INT, VII, p. 56, § 2, no. 5, th. 2, c)).

Let \(p \in [1, +\infty[\) and let \(\pi\) a unitary representation of H in a complex Hilbert space E.

Lemma 6. --- Let \(f \in \mathcal{K}_{\overline{\pi}}(\mathrm{G})\) . For every \(g \in \mathrm{G}\) , the function

\[
\widetilde {f} \colon x \mapsto \eta (g ^ {- 1}, x \mathrm{H}) ^ {1 / p} f (g ^ {- 1} x)
\]

of G in E belongs to \(\mathcal{K}_{\overline{\pi}}(\mathrm{G})\) and satisfies \(\mathrm{N}_p(\widetilde{f}) = \mathrm{N}_p(f)\) .

One verifies without difficulty that \(\tilde{f} \in \mathcal{K}_{\overline{\pi}}(\mathrm{G})\) . Since

\[
\mathrm{N} _ {p} (\widetilde {f}) ^ {p} = \int_ {\mathrm{G/H}} ^ {*} \| \widetilde {f} \| ^ {p} d \nu = \int_ {\mathrm{G/H}} ^ {*} \gamma_ {\mathrm{G/H}} (g) (\| f \| ^ {p}) (y) \eta (g ^ {- 1}, y) d \nu (y)
\]

 and that \((y\mapsto \eta (g^{-1},y))\cdot \nu = \pmb{\gamma}_{\mathrm{G / H}}(g)\nu\) , we obtain \(\mathrm{N}_p(\widetilde{f}) = \mathrm{N}_p(f)\) .

 It follows from this lemma that there exists a continuous and isometric representation \(\widetilde{\pi}\) from G into \(\mathrm{L}_{\overline{\pi}}^{p}(\mathrm{G})\) such that for \(f\in \mathcal{K}_{\pi}(\mathrm{G})\) and \(g\in \mathrm{G}\) , the element \(\widetilde{\pi}(g)f\) is the class of the function \(\widetilde{f}\) defined above. If \(p=2\) , then this representation is unitary.

DEFINITION 2. --- One says that the unitary representation of G in the space \(L_{\overline{\pi}}^{2}(G)\) thus defined is the unitary representation of G induced by the representation \(\pi\) of H with respect to \(\kappa\) . We denote it by \(\operatorname{Ind}_{\mathrm{H}}^{\mathrm{G}}(\pi, \kappa)\) , or simply \(\operatorname{Ind}_{\mathrm{H}}^{\mathrm{G}}(\pi)\) .

 Remarks. --- 1) Let \(\varrho\) a unitary representation of H in a complex Hilbert space F, and let \(u\colon\pi\to\varrho\) an H-morphism. For every function \(f\in\mathcal{K}_{\overline{\pi}}(\mathrm{G})\) , denote \(v(f)\) the function \(g\mapsto u(f(g))\) from G to F; it belongs to \(\mathcal{K}_{\overline{\varrho}}(\mathrm{G})\) and satisfies \(\mathrm{N}_{p}(v(f))\leqslant\|u\|\mathrm{N}_{p}(f)\) . The linear mapping from \(\mathcal{K}_{\overline{\pi}}(\mathrm{G})\) in \(\mathcal{K}_{\overline{\varrho}}(\mathrm{G})\) which to \(f\) associates \(v(f)\) extends therefore to a continuous H-morphism from \(\mathrm{Ind}_{\mathrm{H}}^{\mathrm{G}}(\pi)\) in \(\mathrm{Ind}_{\mathrm{H}}^{\mathrm{G}}(\varrho)\) which is denoted \(\mathrm{Ind}_{\mathrm{H}}^{\mathrm{G}}(u)\) . We have \(\mathrm{Ind}_{\mathrm{H}}^{\mathrm{G}}(1_{\pi})=1_{\mathrm{Ind}_{\mathrm{H}}^{\mathrm{G}}(\pi)}\) . Let \(\sigma\) a unitary representation of H, and let \(v\colon\varrho\to\sigma\) an H-morphism; we have \(\mathrm{Ind}_{\mathrm{H}}^{\mathrm{G}}(v\circ u)=\mathrm{Ind}_{\mathrm{H}}^{\mathrm{G}}(v)\circ\mathrm{Ind}_{\mathrm{H}}^{\mathrm{G}}(u)\) .

In other words, the construction which to \(\pi\) associates \(\mathrm{Ind}_{\mathrm{H}}^{\mathrm{G}}(\pi)\) and with \(u\) associates \(\mathrm{Ind}_{\mathrm{H}}^{\mathrm{G}}(u)\) is a functor from the category of unitary representations of H to that of unitary representations of G (cf. CAT, I, § 2, in preparation).

\begin{enumerate}
\def\labelenumi{\arabic{enumi})}
\setcounter{enumi}{1}
\tightlist
\item
  Let \(\kappa^{\prime}\colon \mathrm{G}\to \mathbf{R}_{+}^{*}\) such that
\end{enumerate}

\[
\frac {\kappa^ {\prime} (x h)}{\kappa^ {\prime} (x)} = \frac {\Delta_ {\mathrm{H}} (h)}{\Delta_ {\mathrm{G}} (h)} = \frac {\kappa (x h)}{\kappa (x)}
\]

 for all \((x,h)\in\mathrm{G}\times\mathrm{H}\) . Let \(\nu^{\prime}\) the quasi-invariant measure \((\kappa^{\prime}\cdot\mu)/\beta\) on G/H. The function \(\kappa^{\prime}\kappa^{-1}\) defines, by passing to the quotient, a continuous function \(\xi\colon G/H\to R_{+}^{*}\) such that \(\nu^{\prime}=\xi\cdot\nu\) . The endomorphism \(\alpha\) of \(\mathcal{K}_{\pi}(G)\) defined by \(f\mapsto(\kappa^{\prime}\kappa^{-1})^{1/p}f\) satisfies

\[
\int_ {\mathrm{G} / \mathrm{H}} ^ {*} \| f \| ^ {p} d \nu^ {\prime} = \int_ {\mathrm{G} / \mathrm{H}} ^ {*} \| \alpha (f) \| ^ {p} d \nu
\]

and makes it possible to identify the spaces \(\mathcal{L}_{\pi}^{p}(\mathrm{G},\nu^{\prime})\) and \(\mathcal{L}_{\pi}^{p}(\mathrm{G},\nu)\) , as well as the spaces \(\mathrm{L}_{\pi}^{p}(\mathrm{G},\nu^{\prime})\) and \(\mathrm{L}_{\pi}^{p}(\mathrm{G},\nu)\) . Moreover, \(\alpha\) induces an isometric isomorphism of representations \(\mathrm{Ind}_{\mathrm{H}}^{\mathrm{G}}(\pi,\kappa^{\prime})\) and \(\mathrm{Ind}_{\mathrm{H}}^{\mathrm{G}}(\pi,\kappa)\) .

\subsection{Case of a closed central subgroup}

In this issue, we assume that the group G is unimodular.

 Consider a closed subgroup Z of the center of G, and let dz denote a Haar measure on Z. We endow the quotient group G/Z with the Haar measure \(\nu = \mu/dz\) . Let \(\chi\) a unitary character of Z. The quotient group G/Z is unimodular by INT, VII, p. 61, § 2, no. 7, cor. of prop. 11.

Lemma 7. --- One has

\[
\mathrm{N} _ {1} (\check {f}) = \mathrm{N} _ {1} (f), \quad \langle f _ {1} | f _ {2} \rangle = \langle \check {f} _ {1} | \check {f} _ {2} \rangle .\tag{6}
\]

for all \(f \in \mathcal{F}_{\chi}(\mathrm{G})\) and for all \(f_1\) and \(f_2\) in \(\mathcal{L}_{\chi}^{2}(\mathrm{G})\) .

These formulas are consequences of the definitions.

 For every \(g \in \mathbf{G}\) and every \(f \in \mathcal{F}_{\chi}(\mathbf{G})\) , the functions \(x \mapsto f(g^{-1}x)\) and \(x \mapsto f(xg)\) belong to \(\mathcal{F}_{\chi}(\mathbf{G})\) . They are denoted respectively \(\boldsymbol{\gamma}_{\mathrm{G},\chi}(g)f\) and \(\boldsymbol{\delta}_{\mathrm{G},\chi}(g)f\) . The maps \(\boldsymbol{\gamma}_{\mathrm{G},\chi}\) and \(\boldsymbol{\delta}_{\mathrm{G},\chi}\) are linear representations of \(\mathbf{G}\) in \(\mathcal{F}_{\chi}(\mathbf{G})\) . For every \(z \in \mathbf{Z}\) , we have

\[
\boldsymbol {\gamma} _ {\mathrm{G}, \chi} (g z) = \overline {{\chi (z)}} \boldsymbol {\gamma} _ {\mathrm{G}, \chi} (g), \quad \boldsymbol {\delta} _ {\mathrm{G}, \chi} (g z) = \chi (z) \boldsymbol {\delta} _ {\mathrm{G}, \chi} (g).\tag{7}
\]

Let \(f\in \mathcal{F}_{\chi}(\mathrm{G})\) and \(g\in \mathrm{G}\) ; we have

\[
| \boldsymbol {\gamma} _ {\mathrm{G}, \chi} (g) f | = \boldsymbol {\gamma} _ {\mathrm{G} / \mathrm{Z}} (g \mathrm{Z}) | f |, \quad | \boldsymbol {\delta} _ {\mathrm{G}, \chi} (g) f | = \boldsymbol {\delta} _ {\mathrm{G} / \mathrm{Z}} (g \mathrm{Z}) | f |,\tag{8}
\]

where all functions appearing in these equalities are identified with functions on G/Z.

 The subspace \(\mathcal{K}_{\chi}(\mathrm{G})\) of \(\mathcal{F}_{\chi}(\mathrm{G})\) is stable under \(\gamma_{\mathrm{G},\chi}\) and \(\delta_{\mathrm{G},\chi}\) . Let \(p\) a real number \(\geqslant 1\) . The formulas (8) imply that the representations \(\gamma_{\mathrm{G},\chi}\) and \(\delta_{\mathrm{G},\chi}\) restricted to \(\mathcal{K}_{\chi}(\mathrm{G})\) , extend to continuous, isometric linear representations of G into \(\mathcal{L}_{\chi}^{p}(\mathrm{G})\) which will be denoted \(\gamma_{\mathrm{G},\chi}^{(p)}\) and \(\delta_{\mathrm{G},\chi}^{(p)}\) . By passing to quotients, these representations also define isometric representations of G in \(\mathrm{L}_{\chi}^{p}(\mathrm{G})\) denoted in the same way.

 The representations \(\gamma_{\mathrm{G},\chi}^{(2)}\) and \(\delta_{\mathrm{G},\chi}^{(2)}\) in \(\mathrm{L}_{\chi}^{2}(\mathrm{G})\) are unitary, and will be denoted simply \(\gamma_{\mathrm{G},\chi}\) and \(\delta_{\mathrm{G},\chi}\) , respectively, when there is no confusion with the representations in \(\mathcal{F}_{\chi}(\mathrm{G})\) . We also denote \(\varrho_{\mathrm{G},\chi}\) the continuous representation of \(\mathrm{G} \times \mathrm{G}\) in \(\mathcal{L}_{\chi}^{2}(\mathrm{G})\) or \(\mathrm{L}_{\chi}^{2}(\mathrm{G})\) defined by

\[
\pmb {\varrho} _ {\mathrm{G}, \chi} (g, h) = \pmb {\gamma} _ {\mathrm{G}, \chi} (g) \circ \pmb {\delta} _ {\mathrm{G}, \chi} (h) = \pmb {\delta} _ {\mathrm{G}, \chi} (h) \circ \pmb {\gamma} _ {\mathrm{G}, \chi} (g)
\]

for all \((g,h)\in \mathrm{G}\times \mathrm{G}\) (cf. Lemma 1 of V, p. 377).

The representation \(\gamma_{\mathrm{G},\chi}\) onto \(\mathrm{L}_{\chi}^{2}(\mathrm{G})\) is none other than the induced representation \(\mathrm{Ind}_{\mathrm{Z}}^{\mathrm{G}}(\overline{\chi})\) .

When \(Z = \{e\}\) the following lemma results from INT, VIII, p. 166, § 4, no. 5, prop. 12, since G is unimodular.

Lemma 8. --- Let \(f_1 \in \mathcal{K}_\chi(\mathrm{G})\) and \(f_2 \in \mathcal{L}_\chi^2(\mathrm{G})\) .

\begin{enumerate}
\def\labelenumi{\alph{enumi})}
\item
  The function \(f\) onto \(\mathrm{G}\) defined by \(f(g) = \langle f_1 | \boldsymbol{\gamma}_{\mathrm{G},\chi}(g)f_2 \rangle\) for all \(g \in \mathrm{G}\) belongs to \(\mathcal{L}_{\chi}^{2}(\mathrm{G})\) and satisfies \(\mathrm{N}_2(f) \leqslant \mathrm{N}_1(f_1)\mathrm{N}_2(f_2)\) ;
\item
  The function \(f\) onto \(G\) defined by \(f(g) = \langle f_1 | \delta_{G,\chi}(g)f_2 \rangle\) for all \(g \in G\) belongs to \(\mathcal{L}_{\chi}^{2}(G)\) and satisfies \(N_2(f) \leqslant N_1(f_1)N_2(f_2)\) .
\end{enumerate}

Let us prove \(a\) ), the proof of assertion \(b\) ) being similar. The function \(f\) is continuous, hence \(\mu\) -measurable. For every \(z \in \mathbf{Z}\) and every \(g \in \mathbf{G}\) , we have

\[
f (g z) = \langle f _ {1} | \gamma_ {\mathrm{G}, \chi} (g z) f _ {2} \rangle = \overline {{\chi (z)}} \langle f _ {1} | \gamma_ {\mathrm{G}, \chi} (g) f _ {2} \rangle
\]

(formula (7), p. 417), hence \(f \in \mathcal{F}_{\overline{\chi}}(\mathrm{G})\) . Prop. 5 of V, p. 413 implies that it now suffices to prove that \(\mathrm{N}_2(f) \leqslant \mathrm{N}_1(f_1)\mathrm{N}_2(f_2)\) .

Suppose first that \(f_2\) belongs to \(\mathcal{K}_{\chi}(\mathrm{G})\) . For every \(g \in \mathrm{G}\) , one has by definition

\[
f (g) = \int_ {\mathrm{G/Z}} \overline {{f _ {1}}} \gamma_ {\mathrm{G}, \chi} (g) f _ {2} d \nu ,
\]

where the function \(\overline{f_1} \gamma_{\mathrm{G},\chi}(g)f_2\) is identified with a function on G/Z.

Let us define a function \(f_3\) onto \(G\) by setting \(f_3(g) = 0\) if \(f_1(g) = 0\) and \(f_3(g) = f_1(g)|f_1(g)|^{-1/2}\) otherwise. The function \(f_3\) belongs to \(\mathcal{F}_{\chi}(G)\) and satisfies \(f_1 = |f_1|^{1/2}f_3\) ; it is \(\mu\) -measurable and zero outside a compact modulo \(Z\) since \(f_1\) is. Since \(|f_1|^{1/2} \in \mathcal{K}_1(G)\) , it follows that

\[
\overline {{f _ {1}}} \gamma_ {\mathrm{G}, \chi} (g) f _ {2} = | f _ {1} | ^ {1 / 2} \overline {{f _ {3}}} \gamma_ {\mathrm{G}, \chi} (g) f _ {2},
\]

\[
| \overline {{f _ {3}}} \gamma_ {\mathrm{G}, \chi} (g) f _ {2} | = | f _ {3} | | \gamma_ {\mathrm{G}, \chi} (g) f _ {2} |
\]

for all \(g \in \mathbf{G}\) .

 Let \(g \in G\) . By the Cauchy--Schwarz inequality, we have

\[
\begin{array}{l} | f (g) | ^ {2} \leqslant \Big (\int_ {\mathrm{G/Z}} | f _ {1} | ^ {1 / 2} | f _ {3} | | \boldsymbol {\gamma} _ {\mathrm{G}, \chi} (g) f _ {2} | d \nu \Big) ^ {2} \\ \leqslant \Big (\int_ {\mathrm{G/Z}} | f _ {1} | d \nu \Big) \Big (\int_ {\mathrm{G/Z}} | f _ {3} | ^ {2} | \boldsymbol {\gamma} _ {\mathrm{G}, \chi} (g) f _ {2} | ^ {2} d \nu \Big). \end{array}
\]

Since we have \(|f_{3}|^{2}=|f_{1}|\) in \(\mathcal{K}(\mathrm{G}/\mathrm{Z})\) , we deduce by integrating over G/Z that

\[
\int_ {\mathrm{G} / \mathrm{Z}} | f (g) | ^ {2} d \nu (g) \leqslant \mathrm{N} _ {1} (f _ {1}) \int_ {\mathrm{G} / \mathrm{Z}} \left(\int_ {\mathrm{G} / \mathrm{Z}} | f _ {1} (x) | | f _ {2} (g ^ {- 1} x) | ^ {2} d \nu (x)\right) d \nu (g).
\]

The function on G/Z × G/Z deduced from the function

\[
(g, x) \mapsto | f _ {1} (x) | | f _ {2} (g ^ {- 1} x) | ^ {2}
\]

by passing to quotients is \((\nu \otimes \nu)\) measurable and compactly supported, hence it is \((\nu \otimes \nu)\) -moderate (INT, V, p. 4, § 5, no. 2, def. 2). By prop. 7 of INT, V, p. 93, § 8, no. 3, it follows that

\[
\begin{array}{l} \int_ {\mathrm{G} / \mathrm{Z}} \Bigl (\int_ {\mathrm{G} / \mathrm{Z}} | f _ {1} (x) |   | f _ {2} (g ^ {- 1} x) | ^ {2} d \nu (x) \Bigr) d \nu (g) \\ = \int_ {\mathrm{G} / \mathrm{Z}} | f _ {1} (x) | \Bigl (\int_ {\mathrm{G} / \mathrm{Z}} | f _ {2} (g ^ {- 1} x) | ^ {2} d \nu (g) \Bigr) d \nu (x) = \mathrm{N} _ {1} (f _ {1}) \mathrm{N} _ {2} (f _ {2}) ^ {2}. \end{array}
\]

Consequently, we have \(\mathrm{N}_{2}(f)^{2}\leqslant\mathrm{N}_{1}(f_{1})^{2}\mathrm{N}_{2}(f_{2})^{2}\) which establishes the desired property in this case.

 Let us consider the general case. Denote by \(u\) the linear map from \(\mathcal{K}_{\chi}(\mathrm{G})\) in \(\mathcal{L}_{\overline{\chi}}^{2}(\mathrm{G})\) which to \(f_{2}\) associates \(f\) . Let \(f_{2} \in \mathcal{L}_{\chi}^{2}(\mathrm{G})\) and let \((f_{2,n})_{n \in \mathbf{N}}\) a sequence in \(\mathcal{K}_{\chi}(\mathrm{G})\) which converges to \(f_{2}\) in \(\mathcal{L}_{\chi}^{2}(\mathrm{G})\) . Let \(f_{n} = u(f_{2,n})\) ; the sequence \((f_{n})_{n \in \mathbf{N}}\) is Cauchy in \(\mathcal{L}_{\overline{\chi}}^{2}(\mathrm{G})\) since the previous case implies that \(\mathrm{N}_{2}(f_{n} - f_{m}) \leqslant \mathrm{N}_{1}(f_{1})\mathrm{N}_{2}(f_{2,n} - f_{2,m})\) for all \(n\) and \(m\) in \(\mathbf{N}\) . Let \(f \in \mathcal{L}_{\overline{\chi}}^{2}(\mathrm{G})\) such that \((f_{n})\) converges to \(f\) (prop. 3, \(c\) ) of V, p. 409). Since \(\mathrm{N}_{2}(f_{n}) \leqslant \mathrm{N}_{1}(f_{1})\mathrm{N}_{2}(f_{2,n})\) for all \(n \in \mathbf{N}\) , it follows that \(\mathrm{N}_{2}(f) \leqslant \mathrm{N}_{1}(f_{1})\mathrm{N}_{2}(f_{2})\) .

There exists a subsequence \((f_{n_{k}})_{k\in\mathbf{N}}\) such that \(f_{n_{k}}(g)\) converges to \(f(g)\) for all \(g\in G\) outside of a negligible subset of G modulo Z (loc. cit., b)). But on the other hand, for every \(g\in G\) , we have

\[
f _ {n _ {k}} (g) = \left\langle f _ {1} \mid \gamma_ {\mathrm{G}, \chi} (g) f _ {2, n _ {k}} \right\rangle\rightarrow \left\langle f _ {1} \mid \gamma_ {\mathrm{G}, \chi} (g) f _ {2} \right\rangle .
\]

Consequently, we have \(f(g) = \langle f_1 | \gamma_{\mathrm{G},\chi}(g)f_2 \rangle\) for all \(g\) outside a negligible subset of G modulo Z. Since \(f \in \mathcal{L}_{\overline{\chi}}^2(\mathrm{G})\) and that \(\mathrm{N}_2(f) \leqslant \mathrm{N}_1(f_1)\mathrm{N}_2(f_2)\) , the lemma is proved.

\subsection{Square-integrable representations}

In this issue, the Hilbert spaces considered are complex. It is assumed that G is unimodular, and its center is denoted by C. For every closed subgroup Z of C, we denote by \(\beta_{Z}\) a Haar measure on Z, and we equip G/Z with the Haar measure \(\nu_{Z} = \mu/\beta_{Z}\) (INT, VII, p. 44, § 2, no. 2, def. 1).

Let \(\pi\) an irreducible unitary representation of G in a Hilbert space E and \(\chi \in \widehat{\mathbf{C}}\) its central character (def. 6 of V, p. 390). For all \(x\) and \(y\) in E, we denote \(f_{x,y}\) the matrix coefficient \(g \mapsto \langle x \mid \pi(g)y \rangle\) ; it is a continuous function on G with complex values.

Let \(x\) and \(y\) in E. For \(z \in \mathrm{C}\) and \(g \in \mathrm{G}\) , we have

\[
f _ {x, y} (z g) = \langle x \mid \pi (z g) y \rangle = \langle x \mid \chi (z) \pi (g) y \rangle = \chi (z) f _ {x, y} (g),\tag{9}
\]

hence \(f_{x,y}\) belongs to the space \(\mathcal{F}_{\chi}(\mathrm{G})\) (V, p. 408). Moreover, for every \((g_1, h_1) \in \mathrm{G} \times \mathrm{G}\) and every \(g \in \mathrm{G}\) , we have

\[
f _ {\pi (g _ {1}) x, \pi (h _ {1}) y} (g) = \langle \pi (g _ {1}) x | \pi (g) \pi (h _ {1}) y \rangle = f _ {x, y} (g _ {1} ^ {- 1} g h _ {1}).\tag{10}
\]

Relation (9) justifies the following definition.

DEFINITION 3. --- Let \(\pi\) an irreducible unitary representation of G in a Hilbert space E. We say that \(\pi\) is square-integrable modulo the center if the function on G/C deduced from the function \(|f_{x,y}|\) by passage to the quotient belongs to \(\mathcal{L}^2 (\mathrm{G / C})\) for all \((x,y)\in \mathrm{E}\times \mathrm{E}\) .

This condition does not depend on the choice of a Haar measure on G/C.

If the matrix coefficients of \(\pi\) belong to \(\mathcal{L}^2 (\mathrm{G})\) , one says that \(\pi\) is square-integrable; the existence of an irreducible unitary representation of G that is square-integrable implies that the center of G is compact (exercise 5 of V, p. 487).

There exist groups G which admit no square-integrable representation, even modulo the center (cf. exercise 32 of V, p. 516).

PROPOSITION 6. --- Let \(\pi\) an irreducible unitary representation of G in a Hilbert space E. Let \(\chi\) the central character of \(\pi\) . Then \(\pi\) is square-integrable modulo the center if and only if there exist nonzero elements \(x_0\) and \(y_0\) of E such that the function on G/C induced by \(|f_{x_0,y_0}|\) by passage to the quotient belongs to \(\mathcal{L}^2 (\mathrm{G / C})\) .

Suppose there exist elements \(x_{0}\) and \(y_{0}\) nonzero elements of E such that \(|f_{x_{0},y_{0}}| \in \mathcal{L}^{2}(\mathrm{G/C})\) . It suffices to show that \(\pi\) is then square-integrable modulo the center.

 Let F be the set of elements \(x \in E\) such that \(|f_{x,y_{0}}|\) belongs to \(\mathcal{L}^{2}(G/C)\) . This is a vector subspace of E; it contains \(x_{0}\) , hence is nonzero. Relation (10) above implies that F is stable under \(\pi\) ; since the representation \(\pi\) is irreducible, the space F is therefore dense in E.

 Let \(x \in \mathrm{F}\) . Since \(f_{x,y_0}\) belongs to \(\mathcal{F}_{\chi}(\mathrm{G})\) and is \(\mu\) -measurable, and \(|f_{x,y_0}| \in \mathcal{L}^2(\mathrm{G/C})\) , we have \(f_{x,y_0} \in \mathcal{L}_{\chi}^2(\mathrm{G})\) (prop. 5 of V, p. 413 applied to Z = C). Let \(u\) denote the partial operator from E to \(\mathrm{L}_{\chi}^{2}(\mathrm{G})\) whose domain is F and which associates to \(x \in \mathrm{F}\) the class of \(f_{x,y_0}\) in \(\mathrm{L}_{\chi}^{2}(\mathrm{G})\) .

 Let us show that the partial operator \(u\) is closed. Let \((x_{n}, u(x_{n}))_{n \in \mathbf{N}}\) a sequence of elements of the graph of \(u\) which converges in \(\mathrm{E} \times \mathrm{L}_{\chi}^{2}(\mathrm{G})\) . Let \(x\) the limit of \((x_{n})\) and let \(f \in \mathcal{L}_{\chi}^{2}(\mathrm{G})\) a function whose class is the limit of the sequence \((u(x_{n}))\) .

The function \(u(x_{n})\) is the class of the matrix coefficient \(f_{x_{n},y_{0}}\) . For every \(g \in G\) , we have

\[
f _ {x _ {n}, y _ {0}} (g) = \langle x _ {n} | \pi (g) y _ {0} \rangle \rightarrow \langle x | \pi (g) y _ {0} \rangle = f _ {x, y _ {0}} (g)
\]

as \(n \to +\infty\) . We therefore have \(f_{x,y_{0}} \in \mathcal{L}_{\chi}^{2}(\mathrm{G})\) and \(f = f_{x_{0},y}\) almost everywhere modulo C (cor. of prop. 3 of V, p. 409); this means that \(u(x)\) is the class of f in \(\mathrm{L}_{\chi}^{2}(\mathrm{G})\) . We have therefore proved that u is closed.

The domain F of u is stable under \(\pi\) , and relation (10) shows that u satisfies \(u \circ \pi(g) = \gamma_{\mathrm{G},\chi}(g) \circ u\) for all \(g \in G\) . Consequently we also have the equality \(u^{*} \circ \gamma_{\mathrm{G},\chi}(g) = \pi(g) \circ u^{*}\) for all \(g \in G\) (lemma 6 of V, p. 381), whence \((u^{*} \circ u) \circ \pi(g) = \pi(g) \circ (u^{*} \circ u)\) . Now the partial operator \(u^{*} \circ u\) is self-adjoint (prop. 12 of IV, p. 241), and in particular closed (prop. 7 of IV, p. 236), so cor. 1 of V, p. 387 implies that the domain of \(u^{*} \circ u\) is equal to E. A fortiori, we have F = E, that is, the function \(|f_{x,y_{0}}|\) belongs to \(\mathcal{L}^{2}(\mathrm{G}/\mathrm{C})\) for all \(x \in E\) .

Let \((x,y)\in\mathrm{E}\times\mathrm{E}\) . We have \(f_{x,y_{0}}\in\mathcal{L}_{\chi}^{2}(\mathrm{G})\) . One proves mutatis mutandis, using the representation \(\delta_{G,\chi}\) instead of \(\gamma_{G,\chi}\) , that the set of \(y\in E\) such that the function \(|f_{x,y}|\) belongs to \(\mathcal{L}^{2}(\mathrm{G}/\mathrm{C})\) is equal to E. The proposition follows.

In the rest of this section, we fix a closed subgroup Z of C such that C/Z is compact.

Lemma 9. --- Let \(\pi\) an irreducible unitary representation of G in a Hilbert space E which is square-integrable modulo the center. Let \(\chi\) the restriction to Z of the central character of \(\pi\) . For all \(x\) and \(y\) in E, we have \(f_{x,y} \in \mathcal{L}_{\chi}^{2}(\mathrm{G})\) .

 The function \(f_{x,y}\) is continuous hence \(\mu\) -measurable. We have \(f_{x,y} \in \mathcal{F}_{\chi}(\mathrm{G})\) by formula (9), p. 420. Moreover, by INT, VII, p. 64, § 2, no. 8, cor. 1, c), we have \(\mathrm{N}_{2}(f_{x,y}) < +\infty\) since C/Z is compact and \(|f_{x,y}| \in \mathcal{L}^{2}(\mathrm{G}/\mathrm{C})\) . The assertion therefore follows from Proposition 5 of V, p. 413.

PROPOSITION 7. --- Let π be an irreducible unitary representation of G in a Hilbert space E which is square-integrable modulo the center. Let χ be the restriction to Z of the central character of π.

There exists a real number c \textgreater{} 0 and a unique \((\mathbf{G} \times \mathbf{G})\) -isometric morphism w of the unitary representation \(\overline{\pi} \boxtimes \pi\) in \(\mathrm{L}_{\chi}^{2}(\mathrm{G})\) such that, for all \((x, y) \in \overline{\mathrm{E}} \times \mathrm{E}\) , the element \(w(x \otimes y)\) is the class in \(\mathrm{L}_{\chi}^{2}(\mathrm{G})\) of the function \(c^{1/2} f_{x,y}\) .

 For every \((x,y)\in\mathrm{E}\times\mathrm{E}\) , we have \(f_{x,y}\in\mathcal{L}_{\chi}^{2}(\mathrm{G})\) (Lemma 9). Let v denote the unique linear map from \(\overline{E}\otimes E\) in \(\mathrm{L}_{\chi}^{2}(\mathrm{G})\) such that \(v(x\otimes y)\) is the class of \(f_{x,y}\) for all \((x,y)\in\overline{\mathrm{E}}\times\mathrm{E}\) .

We prove below the following lemma.

Lemma 10. --- There exists a real number c \textgreater{} 0 such that the linear map \(w = c^{1/2}v\) is isometric.

This lemma being assumed valid, note that formula (10), p. 420, can be written as

\[
v (\overline {{\pi}} (g _ {1}) x \otimes \pi (h _ {1}) y) = \varrho_ {\mathrm{G}, \chi} (g _ {1}, h _ {1}) v (x \otimes y)\tag{11}
\]

 for all \((g_{1}, h_{1}) \in \mathbf{G} \times \mathbf{G}\) and every \((x, y) \in \overline{\mathbf{E}} \times \mathbf{E}\) . The isometric linear map w from \(\overline{\mathbf{E}} \otimes \mathbf{E}\) in \(\mathrm{L}_{\chi}^{2}(\mathrm{G})\) admits a continuous extension, still denoted w, to \(\overline{E} \widehat{\otimes}_{2} E\) . By continuity and linearity, formula (11) implies that w is a \((\mathrm{G} \times \mathrm{G})\) -morphism of \(\overline{\pi} \boxtimes \pi\) in \(\mathrm{L}_{\chi}^{2}(\mathrm{G})\) , which concludes the proof of the proposition.

 Let us prove the lemma. For all \(x \in E\) , we denote \(u_{x}\) the linear map \(y \mapsto v(x \otimes y) = f_{x,y}\) from E into \(\mathrm{L}_{\chi}^{2}(\mathrm{G})\) . We have \(u_{x} \in \mathrm{Hom}_{\mathrm{G}}(\pi, \gamma_{\mathrm{G},\chi})\) . By formula (10), according to Corollary 5 of V, p. 388, there exists a real number \(\lambda_{x} \geqslant 0\) such that \(\lambda_{x}u_{x}\) is isometric.

Let \(x\) and \(y\) in E. We have

\[
\| f _ {x, y} \| ^ {2} = \int_ {\mathrm{G/Z}} \overline {{f}} _ {x, y} f _ {x, y} d \nu_ {\mathrm{Z}} = \int_ {\mathrm{G/Z}} \check {\overline {{f}}} _ {x, y} \check {f} _ {x, y} d \nu_ {\mathrm{Z}}
\]

since G/Z is unimodular (Lemma 7 of V, p. 417). Noting that \(\check{f}_{x,y} = \overline{f}_{y,x}\) , we obtain

\[
\lambda_ {x} \| y \| ^ {2} = \| f _ {x, y} \| ^ {2} = \int_ {\mathrm{G/Z}} f _ {y, x} \overline {{f}} _ {y, x} d \nu_ {\mathrm{Z}} = \| f _ {y, x} \| ^ {2} = \lambda_ {y} \| x \| ^ {2}.
\]

This means that the positive real number \(\lambda_{x}/\|x\|^{2}\) is independent of the choice of the nonzero element x of E. It is strictly positive since, for every nonzero x, the function \(f_{x,x}\) is continuous and takes the value \(\|x\|^{2}>0\) at e, whence \(\|u_{x}(x)\|=\|f_{x,x}\|>0\) . Denote by \(c^{-1}\) this real number.

For every \((x,y)\in\mathrm{E}\times\mathrm{E}\) , it follows that

\[
\| v (x \otimes y) \| ^ {2} = \| f _ {x, y} \| ^ {2} = \lambda_ {x} \| y \| ^ {2} = c ^ {- 1} \| x \| ^ {2} \| y \| ^ {2} = c ^ {- 1} \| x \otimes y \| ^ {2}.
\]

Using EVT, V, p. 29, cor. 1, we deduce that the linear map \(w = c^{1/2}v\) of \(\overline{\mathrm{E}}\widehat{\otimes}_2\mathrm{E}\) in \(\mathrm{L}_{\chi}^{2}(\mathrm{G})\) is isometric, as desired.

 COROLLARY. --- Let \(\pi\) an irreducible unitary representation of G and let \(\chi\) the restriction to Z of its central character. The representation \(\pi\) is square-integrable modulo the center if and only if it is isomorphic to a subrepresentation of the representation \(\gamma_{\mathrm{G},\overline{\chi}}\) from G into \(\mathrm{L}_{\overline{\chi}}^{2}(\mathrm{G})\) (resp. of the representation \(\delta_{\mathrm{G},\chi}\) from G into \(\mathrm{L}_{\chi}^{2}(\mathrm{G})\) ).

Let us prove the assertion concerning \(\gamma_{\mathrm{G},\overline{\chi}}\) , the second being proved in a similar manner.

 Suppose first that there exists a subrepresentation E of the representation \(\gamma_{\mathrm{G},\overline{\chi}}\) isomorphic to \(\pi\) . Since the space E is not zero, there exists a function \(f_1 \in \mathcal{K}_{\overline{\chi}}(\mathrm{G})\) whose class \(\widetilde{f}_1\) that is not orthogonal to E. We have \(f_1 \in \mathcal{L}_{\overline{\chi}}^1 (\mathrm{G})\) . Denote by \(\widetilde{f}_{1,\mathrm{E}}\) the orthogonal projection of \(\widetilde{f}_1\) on E; it is a nonzero element of E. Let moreover \(\widetilde{f}_2 \in \mathrm{E}\) non-zero. The map \(h: g \mapsto \langle \widetilde{f}_{1,\mathrm{E}} | \gamma_{\mathrm{G},\overline{\chi}}(g)\widetilde{f}_2 \rangle\) is a matrix coefficient of \(\pi\) . For every \(g \in \mathrm{G}\) , we have \(h(g) = \langle \widetilde{f}_1 | \gamma_{\mathrm{G},\overline{\chi}}(g)\widetilde{f}_2 \rangle\) since \(\widetilde{f}_1 - \widetilde{f}_{1,\mathrm{E}}\) is orthogonal to E. By Lemma 8 of V, p. 418, the function \(h\) belongs to \(\mathcal{L}_{\chi}^{2}(\mathrm{G})\) , so Proposition 6 implies that \(\pi\) is square-integrable modulo the center.

Conversely, suppose that \(\pi\) is square-integrable modulo the center. Let \(x_0 \in \mathrm{E}\) a nonzero vector. By prop. 7 and formula (10), the map \(y \mapsto \check{f}_{x_0,y}\) is an injective G-morphism from \(\pi\) in \(\gamma_{\mathrm{G},\overline{\chi}}\) .

DEFINITION 4. --- Let Z be a closed subgroup of C such that C/Z is compact. Let π be an irreducible unitary representation of G which is square-integrable modulo the center. The unique real number c \textgreater{} 0 satisfying the property of Proposition 7 is called the formal degree of π relative to Z. It is denoted by \(c_{\mathrm{Z}}(\pi)\) .

The formal degree depends on the choice of Haar measure on Z. If the Haar measure \(\beta_{Z}\) on Z is multiplied by a real number t \textgreater{} 0, then the measure \(\nu_{Z} = \mu/\beta_{Z}\) on G/Z is multiplied by \(t^{-1}\) , and the formal degree of \(\pi\) is multiplied by t.

The formal degree is characterized by the following property:

PROPOSITION 8 (Orthogonality relations). --- Let π be an irreducible unitary representation of G in a Hilbert space E that is square-integrable modulo the center. Then one has

\[
c _ {\mathrm{Z}} (\pi) \int_ {\mathrm{G} / \mathrm{Z}} \overline {{\langle x | \pi (g) x ^ {\prime} \rangle}} \langle y | \pi (g) y ^ {\prime} \rangle d \nu_ {\mathrm{Z}} (g) = \overline {{\langle x | y \rangle}} \langle x ^ {\prime} | y ^ {\prime} \rangle
\]

for all \((x,y,x',y')\in \mathrm{E}^4\)

Denote by \(w\) the morphism of Prop. 7. We have

\[
\int_ {\mathrm{G} / \mathrm{Z}} \overline {{\langle x | \pi (g) x ^ {\prime} \rangle}} \langle y | \pi (g) y ^ {\prime} \rangle d \nu_ {\mathrm{Z}} (g) = \langle f _ {x, x ^ {\prime}} | f _ {y, y ^ {\prime}} \rangle
\]

and according to loc. cit., it follows that

\[
\begin{array}{c} \langle f _ {x, x ^ {\prime}} | f _ {y, y ^ {\prime}} \rangle = \frac {1}{c _ {\mathrm{Z}} (\pi)} \langle w (x \otimes x ^ {\prime}) | w (y \otimes y ^ {\prime}) \rangle \\ = \frac {1}{c _ {\mathrm{Z}} (\pi)} \langle x \otimes x ^ {\prime} | y \otimes y ^ {\prime} \rangle = \frac {1}{c _ {\mathrm{Z}} (\pi)} \langle x | y \rangle_ {\overline {{\mathrm{E}}}} \langle x ^ {\prime} | y ^ {\prime} \rangle_ {\mathrm{E}}, \end{array}
\]

whence the result.

In addition to the previous proposition, we also have the following relations for non-isomorphic square-integrable irreducible representations.

PROPOSITION 9. --- Let \(\pi_1\) and \(\pi_2\) non-isomorphic irreducible unitary representations of G in Hilbert spaces E \(_1\) and E \(_2\) . One assumes that \(\pi_1\) and \(\pi_2\) are square-integrable modulo the center and the restrictions to Z of their central characters coincide. Then

\[
\int_ {\mathrm{G} / \mathrm{Z}} \overline {{\langle x \mid \pi_ {1} (g) x ^ {\prime} \rangle}} \langle y \mid \pi_ {2} (g) y ^ {\prime} \rangle d \nu_ {\mathrm{Z}} (g) = 0
\]

for all \((x,x^{\prime},y,y^{\prime})\in \mathrm{E}_{1}^{2}\times \mathrm{E}_{2}^{2}\)

 For \(i=1,2\) , denote \(w_{i}\) the morphism from Proposition 7 for the representation \(\pi_{i}\) . By Lemma 8, \(b\) ) of V, p. 384 and assertion \(b\) ) of Proposition 8 of V, p. 394, the image of \(w_{i}\) is contained in the component \(\pi_{i}\) -isotypic of \(\delta_{\mathrm{G},\chi}\) . By assertion \(a\) ) of loc. cit., the image of \(w_{1}\) is therefore orthogonal to the image of \(w_{2}\) . Consequently, it follows that \(\langle w_{1}(x\otimes x^{\prime})|w_{2}(y\otimes y^{\prime})\rangle=0\) for all \((x,x^{\prime},y,y^{\prime})\in\mathrm{E}_{1}^{2}\times\mathrm{E}_{2}^{2}\) , which is the desired formula.

Remark. --- The orthogonality relations of Props. 8 and 9 generalize those of A, VIII, p. 399 (see also the case where G is compact in § 2 of V, p. 457).

\subsection{Subrepresentations of the regular representation of a commutative group}

In this section, G is a locally compact commutative group and \(\mu\) a Haar measure on G. We denote by \(\hat{G}\) the dual group of G (def. 2 of II, p. 201) and by \(\hat{\mu}\) the dual Haar measure of \(\mu\) onto \(\hat{G}\) (def. 4 of II, p. 214). The notions of measurability will always be relative to \(\mu\) and \(\hat{\mu}\) .

We propose to determine all subrepresentations of the left regular representation \(\gamma_{G}\) of G in \(L^{2}(G,\mu)\) . Since G is commutative, we also have \(\delta_{\mathrm{G}}(g)=\gamma_{\mathrm{G}}(g^{-1})\) , so these subrepresentations are also the subrepresentations of the right regular representation.

For every measurable subset M of \(\widehat{\mathbf{G}}\) , we denote \(\mathrm{E_M}\) the set of \(f \in \mathrm{L}^2 (\mathrm{G},\mu)\) such that the Fourier transform \(\mathcal{F}_{\mathrm{G}}(f)\) is zero almost everywhere on M (cf. no. 3 of II, p. 210). This is the kernel of the continuous linear map \(f \mapsto \varphi_{\mathrm{M}}\mathcal{F}_{\mathrm{G}}(f)\) of \(\mathrm{L}^2 (\mathrm{G},\mu)\) in \(\mathrm{L}^2 (\widehat{\mathrm{G}},\widehat{\mu})\) , where \(\varphi_{\mathrm{M}}\) denotes the characteristic function of M (cf. th. 1 of II, p. 215), and it is therefore a closed subspace of \(\mathrm{L}^2 (\mathrm{G},\mu)\) .

We say that measurable subsets M and N of \(\widehat{G}\) are equal up to a locally negligible set if \((\mathrm{M} \cup \mathrm{N}) - (\mathrm{M} \cap \mathrm{N})\) is locally negligible. Equivalently, measurable subsets M and N are equal up to a locally negligible set if and only if the characteristic functions of M and N are equal in \(\mathrm{L}^{\infty}(\widehat{\mathrm{G}}, \widehat{\mu})\) .

PROPOSITION 10. --- a) Let M be a measurable subset of \(\widehat{G}\) . The space \(\mathrm{E}_{\mathrm{M}}\) is a subrepresentation of the representation \(\gamma_{\mathrm{G}}\) ;

\begin{enumerate}
\def\labelenumi{\alph{enumi})}
\setcounter{enumi}{1}
\item
  Let M and N be measurable subsets of \(\widehat{\mathbf{G}}\) . We have \(\mathrm{E_M} = \mathrm{E_N}\) if and only if M and N are equal up to a locally negligible set;
\item
  Any subrepresentation of \(\gamma_{\mathrm{G}}\) is of the form \(\mathrm{E}_{\mathrm{M}}\) for a measurable subset M of \(\widehat{\mathrm{G}}\) .
\end{enumerate}

Let \(\eta\) the canonical map from G to \(\widehat{\mathbf{G}}\) (cf. II, p. 216, remark 1). Since \(\mathrm{E_M}\) is closed, the assertion \(a\) ) follows from the formulas

\[
\boldsymbol {\gamma} _ {\mathrm{G}} (x) (f) = \varepsilon_ {x} * f, \quad \mathcal {F} _ {\mathrm{G}} (\varepsilon_ {x} * f) = \eta (x) \mathcal {F} _ {\mathrm{G}} (f)
\]

for \(x\in \mathrm{G}\) and \(f\in \mathrm{L}^2 (\mathrm{G},\mu)\) .

Let M and N be measurable subsets equal up to a locally negligible set. Since \(\varphi_{M}\) is then equal to \(\varphi_{N}\) in \(L^{\infty}(\widehat{G},\widehat{\mu})\) , this condition implies that \(E_{M}=E_{N}\) .

Conversely, suppose that M and N are not equal up to a locally negligible set. Possibly after interchanging the roles of M and N, there then exists a compact subset K such that the set \(L = K \cap (M - (M \cap N))\) is not negligible. Let \(\varphi_{\mathrm{L}} \in \mathrm{L}^{2}(\widehat{\mathrm{G}}, \widehat{\mu})\) the class of the characteristic function of L, and set \(f = \overline{\mathcal{F}}_{\widehat{\mathrm{G}}}(\varphi_{\mathrm{L}}) \in \mathrm{L}^{2}(\mathrm{G}, \mu)\) ; we have \(\mathcal{F}_{\mathrm{G}}(f) = \varphi_{\mathrm{L}}\) (corollary of Theorem 2 in II, p. 220). Then f belongs to \(E_{N}\) since \(L \cap N\) is empty, but not in \(E_{M}\), since \(\varphi_{\mathrm{M}}\mathcal{F}(f) = \varphi_{\mathrm{M}}\varphi_{\mathrm{L}} = \varphi_{\mathrm{L}}\) . We therefore have \(E_{M} \neq E_{N}\) , which proves b).

Let now E be a subrepresentation of \(\gamma_{\mathrm{G}}\) . Let \(p_{\mathrm{E}}\) the orthogonal projector from \(\mathrm{L}^2 (\mathrm{G},\mu)\) with image E, and let \(q_{\mathrm{E}} = \mathcal{F}_{\mathrm{G}}\circ p_{\mathrm{E}}\circ \mathcal{F}_{\mathrm{G}}^{-1}\) . The projector \(p_{\mathrm{E}}\) belongs to \(\mathrm{Hom}_{\mathrm{G}}(\gamma_{\mathrm{G}},\gamma_{\mathrm{G}})\) (prop. 4 of V, p. 383), hence it commutes with \(\gamma_{\mathrm{G}}(f)\) for all \(f\in \mathrm{L}^{1}(\mathrm{G},\mu)\) (cf. V, p. 401). This means that it commutes with the endomorphisms \(\varphi \mapsto f*\varphi\) for \(f\in \mathrm{L}^{1}(\mathrm{G},\mu)\) . By lemma 4 of V, p. 407, consequently, the endomorphism \(q_{\mathrm{E}}\) of \(\mathrm{L}^2 (\widehat{\mathrm{G}},\widehat{\mu})\) commutes with the endomorphism of multiplication by \(g\) for every function \(g\in \mathcal{C}_0(\widehat{\mathrm{G}})\) belonging to the image of the Fourier transform of \(\mathrm{L}^1 (\mathrm{G},\mu)\) in \(\mathcal{C}_0(\widehat{\mathrm{G}})\) (prop. 14 of II, p. 223). Since the image of the Fourier transform is dense in \(\mathcal{C}_0(\widehat{\mathrm{G}})\) (corollary of Prop. 5 of II, p. 209), the continuity of the morphism \(g\mapsto m_g\) implies that \(q_{\mathrm{E}}\) commutes with \(m_g\) for every function \(g\in \mathcal{C}_0(\widehat{\mathrm{G}})\) .

By the corollary of Proposition 7 of IV, p. 188, there therefore exists a function \(\varphi \in \mathcal{L}^{\infty}(\widehat{\mathrm{G}},\widehat{\mu})\) such that \(q_{\mathrm{E}} = m_{\varphi}\) . Since \(q_{\mathrm{E}} = q_{\mathrm{E}}^{2}\) , we have \(m_{\varphi} = m_{\varphi}^{2} = m_{\varphi^{2}}\) , whence \(\varphi = \varphi^{2}\) in \(\mathrm{L}^{\infty}(\widehat{\mathrm{G}},\widehat{\mu})\) (prop. 5 of IV, p. 186); this means that \(\varphi\) is equal in \(\mathrm{L}^{\infty}(\widehat{\mathrm{G}},\widehat{\mu})\) to the class of the characteristic function of a measurable subset N of \(\widehat{\mathrm{G}}\) . Set \(\mathrm{M} = \widehat{\mathrm{G}} - \mathrm{N}\) . Let \(f \in \mathrm{L}^{2}(\mathrm{G},\mu)\) . We have \(f \in \mathrm{E}\) if and only if \(p_{\mathrm{E}}(f) = f\) if and only if \(\varphi \mathcal{F}_{\mathrm{G}}(f) = \mathcal{F}_{\mathrm{G}}(f)\) , which is equivalent to \(f \in \mathrm{E}_{\mathrm{M}}\) .

Note. --- Let \(\chi\) a character of G. If G is not compact, the \(\chi\) -isotypic component of the regular representation of G in L \(^{2}\) (G)

is trivial. If G is compact, the \(\chi\) -isotypic component of the regular representation of G has dimension 1, and the function \(\chi \in \mathrm{L}^{2}(\mathrm{G})\) is a basis for it.

\subsection{Unitary representations of the group R}

In this section, E denotes a complex Hilbert space. The group R is equipped with Lebesgue measure.

Lemma 11. --- Let u be a partial self-adjoint operator on E. For \(t \in \mathbf{R}\) set \(\varrho(t) = e^{itu} \in \mathcal{L}(\mathrm{E})\) . Then the map \(\varrho\) is a unitary representation of \(\mathbf{R}\) .

The operator \(\varrho(t)\) is defined by the universally measurable functional calculus (def. 5 of IV, p. 272); it is an endomorphism of E since the function \(x \mapsto e^{itx}\) is bounded on \(\mathrm{Sp}(u)\) (prop. 5, \(a\) ) of IV, p. 275).

According to prop. 5 of IV, p. 275, we have \(\varrho(0) = 1_{\mathrm{E}}\) , \(\varrho(t)^* = \varrho(-t)\) for all \(t \in \mathbf{R}\) and \(\varrho(t_1 + t_2) = \varrho(t_1)\varrho(t_2)\) for all \(t_1\) and \(t_2\) in \(\mathbf{R}\) . In particular, the endomorphism \(\varrho(t)\) is unitary for every \(t \in \mathbf{R}\) . For every \(x \in \mathrm{E}\) , the map from \(\mathbf{R}\) to E defined by \(x \mapsto \varrho(t)x\) is continuous at \(t = 0\) by prop. 6 of IV, p. 276, hence \(\varrho\) is a unitary representation of \(\mathbf{R}\) in E (V, p. 380, lemma 4).

Lemma 12. --- Let \(\varrho\) a unitary representation of \(\mathbf{R}\) in E. Let D be the set of elements \(x\) of E such that the map \(\psi_x: t \mapsto \varrho(t)x\) is differentiable at 0. The map \(u\) of D into E given by \(x \mapsto i^{-1} \psi_x'(0)\) defines a symmetric partial operator on E.

Let \(f \in \mathcal{D}(\mathbf{R})\) and \(x \in \mathrm{E}\) . Set \(y = \varrho(f)x\) . For every \(h \in \mathbf{R}\) , we have

\[
\begin{array}{r l} & {\psi_ {y} (h) = \varrho (h) \varrho (f) x = \varrho (f) \varrho (h) x} \\ & {\qquad = \int_ {\mathbf {R}} f (t) \varrho (t + h) x d t = \int_ {\mathbf {R}} f (t - h) \psi_ {x} (t) d t.} \end{array}
\]

The function of \(R^{2}\) to C defined by \((t,h)\mapsto f(t-h)\) is infinitely differentiable; its derivative with respect to h is the function defined by \((t,h)\mapsto -f'(t-h)\) , which is bounded by \(\|f'\|_{\infty}\) . When h=0, this derivative is zero for t outside a compact set. Since \(\|\psi_{x}(t)\|=\|x\|\) for all t, it follows from prop. 2 of IV, p. 197 that the function \(\psi_{y}\) is differentiable at 0 and that its derivative is given by

\[
\psi_ {y} ^ {\prime} (0) = - \int_ {\mathbf {R}} f ^ {\prime} (t) \psi_ {x} (t) d t = - \varrho (f ^ {\prime}) x.
\]

We therefore have \(y \in D\). Since the space \(\mathcal{D}(\mathbf{R})\) is dense in \(\mathrm{L}^{1}(\mathbf{R})\), it follows from INT, VIII, p. 139, § 2, no. 7, prop. 10, together with (prop. 4 of IV, p. 202), that the space D is dense in E.

Let \(x_{1}\) and \(x_{2}\) in D. We compute

\[
\begin{array}{c} \langle u (x _ {1}) \mid x _ {2} \rangle = \lim _ {h \to 0} \frac {i}{h} \langle (\varrho (h) - 1 _ {\mathrm{E}}) x _ {1} \mid x _ {2} \rangle \\ = \lim _ {h \to 0} \frac {i}{h} \langle x _ {1} \mid (\varrho (- h) - 1 _ {\mathrm{E}}) x _ {2} \rangle = \langle x _ {1} \mid u (x _ {2}) \rangle . \end{array}
\]

Therefore, the partial operator \(u\) is symmetric.

DEFINITION 5. --- Let \(\varrho\) a unitary representation of \(\mathbf{R}\) in E. The symmetric partial operator defined in Lemma 12 is called the infinitesimal generator of \(\varrho\) .

THEOREM 1 (Stone). --- The mapping σ which to a unitary representation ρ of R in E associates its infinitesimal generator u defines a bijection from the set of unitary representations of R in E onto the set of partial self-adjoint operators on E. The inverse bijection τ associates to a partial self-adjoint operator the unitary representation \(t \mapsto e^{itu}\) .

Let us first prove a few lemmas.

Lemma 13. --- Let \(\varrho\) a unitary representation of \(\mathbf{R}\) in E and let \(u\) its infinitesimal generator.

\begin{enumerate}
\def\labelenumi{\alph{enumi})}
\item
  The domain of u is the set of \(x \in E\) such that the mapping \(\psi_{x} \colon t \mapsto \varrho(t)x\) is differentiable on R;
\item
  For every \(t \in \mathbf{R}\) and every \(x \in \mathrm{dom}(u)\) , we have \(\varrho(t)x \in \mathrm{dom}(u)\) and \(\psi_x'(t) = i u(\varrho(t)x)\) ;
\item
  The partial operator u is essentially self-adjoint.
\end{enumerate}

For every \(x \in \mathrm{E}\) , denote \(\psi_x\) the map from \(\mathbf{R}\) in \(\mathrm{E}\) defined by \(\psi_x(t) = \varrho(t)x\) . For every \(t \in \mathbf{R}\) and every \(h \in \mathbf{R}\) , we have

\[
\psi_ {x} (t + h) - \psi_ {x} (t) = \varrho (t) \big (\psi_ {x} (h) - \psi_ {x} (0) \big),
\]

which proves that \(\operatorname{dom}(u)\) is the space of elements \(x \in E\) such that \(\psi_{x}\) is differentiable on R and establishes that \(\psi_{x}^{\prime}(t) = \varrho(t)\psi_{x}^{\prime}(0) = i\varrho(t)(u(x))\) for all \(t \in R\) .

Let \(x \in E\) and \(t \in R\) . We have \(\psi_{\varrho(t)x}(s) = \psi_x(s + t)\) for all \(s \in R\) . Consequently, one has \(\varrho(t)x \in \text{dom}(u)\) if \(x \in \text{dom}(u)\) , and moreover \(u(\varrho(t)x) = i^{-1}\psi_x'(t) = \varrho(t)u(x)\) by a). We then obtain assertion b).

Let us prove c). Let \(x \in \text{Ker}(u^{*} - i1_{\text{E}})\) . Let us prove that x = 0. Let \(y \in \text{dom}(u)\) , and let f be the function on R defined by \(f(t) = \langle \psi_{y}(t) | x \rangle\) for \(t \in R\) . The function f is bounded since \(\| \psi_{y}(t) \| = \| y \|\) for all \(t \in R\) ; it is differentiable on R and, for every \(t \in R\) , assertion b) implies

\[
\begin{array}{r l} f ^ {\prime} (t) = \langle \psi_ {y} ^ {\prime} (t) | x \rangle & = \langle i u (\varrho (t) y) | x \rangle \\ & = - i \langle \varrho (t) y | u ^ {*} (x) \rangle = - i \langle \psi_ {y} (t) | i x \rangle = f (t) \end{array}
\]

since \(u^{*}(x)=ix\) . We therefore have \(f(t)=f(0)e^{t}\) for all \(t\in R\) (FVR, IV, p. 27). Since f is bounded, the function f is zero and, in particular, we have \(\langle y|x\rangle=f(0)=0\) . Since the space dom(u) is dense in E, we conclude that x=0. Similarly, we show that \(\operatorname{Ker}(u^{*}+i1_{\mathrm{E}})\) is reduced to 0. By Cor. 3 of IV, p. 261, the partial operator u is therefore essentially self-adjoint.

Lemma 14. --- Let u be a partial self-adjoint operator on E, and let \(\varrho(t) = e^{itu}\) the unitary representation of R defined by u. The infinitesimal generator of \(\varrho\) is equal to u.

Let \(x \in \text{dom}(u)\) . Set \(\psi_{x}(t) = \varrho(t)x\) for all \(t \in R\) . For every nonzero real number h, we have

\[
\frac {1}{h} \left(\psi_ {x} (t + h) - \psi_ {x} (t)\right) = \left(\frac {1}{h} \left(e ^ {i h u} - 1 _ {\mathrm{E}}\right)\right) e ^ {i t u} x = \left(\frac {1}{h} \left(e ^ {i h u} - 1 _ {\mathrm{E}}\right)\right) \varrho (t) x.
\]

As h tends to 0, we have

\[
\frac {1}{h} (e ^ {i h t} - 1) \rightarrow i t
\]

for all \(t \in R\) . Moreover,

\[
\left| \frac {1}{h} (e ^ {i h t} - 1) \right| = | t | \left| \frac {1}{h} \int_ {0} ^ {h} e ^ {i t s} d s \right| \leqslant | t |.
\]

It then follows from prop. 6 of IV, p. 276 that the function \(\psi_x\) is differentiable on \(\mathbf{R}\) and satisfies \(\psi_x'(t) = i u(\varrho(t)x)\) for all \(t \in \mathbf{R}\) . Consequently, the domain of \(u\) is contained in the set of \(x \in E\) such that \(\psi_x\) is differentiable at 0 and that we then have \(\psi_x'(0) = i u(x)\) . This means by definition that the infinitesimal generator of \(\varrho\) is an extension of \(u\) . These two operators are therefore equal since they are symmetric and \(u\) is self-adjoint (IV, p. 238, remark 5).

Lemma 15. --- Let \(\varrho\) a unitary representation of \(\mathbf{R}\) in E. Then the infinitesimal generator \(u\) of \(\varrho\) is self-adjoint and \(\varrho(t) = e^{itu}\) for all \(t \in \mathbf{R}\) .

By Lemma 13(c), the partial operator u is essentially self-adjoint. Its closure \(\overline{u}\) is therefore a self-adjoint operator. Denote \(\pi\) the unitary representation of R defined by \(\pi(t) = e^{it\overline{u}}\) (lemma 11).

For every \(x \in \mathrm{E}\) , denote \(\psi_x\) (resp. \(\widetilde{\psi}_x\) ) the map from \(\mathbf{R}\) in \(\mathrm{E}\) defined by \(\psi_x(t) = \varrho(t)x\) (respectively by \(\widetilde{\psi}_x(t) = \pi(t)x\) ).

By Lemma 13, parts (a) and (b), the space \(\mathrm{dom}(u)\) is the subspace of E consisting of the elements \(x \in \mathrm{E}\) such that the mapping \(\psi_x\) is differentiable on \(\mathbf{R}\) , and for all \(x \in \mathrm{dom}(u)\) and every \(t \in \mathbf{R}\) , we have \(\psi_x'(t) = i u (\psi_x(t))\) . Likewise, the space \(\mathrm{dom}(\overline{u})\) is the subspace of E consisting of the elements \(x \in \mathrm{E}\) such that the mapping \(\widetilde{\psi}_x\) is differentiable on \(\mathbf{R}\) , and for all \(x \in \mathrm{dom}(\overline{u})\) and every \(t \in \mathbf{R}\) , we have \(\widetilde{\psi}_x'(t) = i \overline{u} (\widetilde{\psi}_x(t))\) .

Let \(x \in \text{dom}(u) \subset \text{dom}(\overline{u})\) . Set \(f = \psi_x - \widetilde{\psi}_x\) . This is a differentiable function from \(\mathbf{R}\) into E. For every \(t \in \mathbf{R}\) , it follows that

\[
f ^ {\prime} (t) = i u (\psi_ {x} (t)) - i \overline {{u}} (\widetilde {\psi} _ {x} (t)) = i \overline {{u}} (f (t))
\]

since \(\psi_x(t) \in \mathrm{dom}(u)\) and \(u \subset \overline{u}\) . Set \(g = \|f\|^2\) ; it is a differentiable map from \(\mathbf{R}\) in \(\mathbf{R}\) such that \(g(0) = 0\) . For \(t \in \mathbf{R}\) , we obtain by FVR, I, p. 28, prop. 2

\[
\begin{array}{r l} g ^ {\prime} (t) & = \langle f ^ {\prime} (t) | f (t) \rangle + \langle f (t) | f ^ {\prime} (t) \rangle \\ & = \langle i \overline {{u}} (f (t)) | f (t) \rangle + \langle f (t) | i \overline {{u}} (f (t)) \rangle = 0 \end{array}
\]

since \(\overline{u}\) is self-adjoint. Hence f = 0, so \(\varrho(t)x = \pi(t)x\) for all \(t \in R\) . The continuous endomorphisms \(\varrho(t)\) and \(\pi(t)\) of E coincide on dom(u), and are therefore equal for all \(t \in R\) . Thus, we have \(\pi = \varrho\) ; since \(\overline{u}\) is the infinitesimal generator of \(\pi\) (Lemma 14), we have \(\overline{u} = u\) , which proves that u is self-adjoint.

We can now prove Theorem 1.

The maps \(\sigma\) and \(\tau\) are well-defined (Lemma 15 and Lemma 11, respectively).

Let \(\varrho\) a unitary representation of \(\mathbf{R}\) . Denote by \(u\) its infinitesimal generator. The relation \(\varrho(t) = e^{itu}\) for all \(t \in \mathbf{R}\) (lemma 15) proves that \(\tau \circ \sigma\) is the identity map.

Let u be a partial self-adjoint operator on E. Lemma 14 shows that the infinitesimal generator of the unitary representation \(t \mapsto e^{itu}\) is equal to u, hence \(\sigma \circ \tau\) is the identity map.

Let u be a partial self-adjoint operator on E and let \(\varrho(t) = e^{itu}\) for \(t \in R\) the unitary representation of R in E associated with it.

Let \(x \in \text{dom}(u)\) . The equation \(\partial_{t} \varrho(t)x = i u(\varrho(t)x)\) which is then satisfied (Lemma 13, b)) is called the Schrödinger equation.

COROLLARY. --- Let \(\varrho\) a unitary representation of \(\mathbf{R}\) in a Hilbert space E. There exists a locally compact space X, a positive measure \(\mu\) on X and a continuous function \(g\) on X with real values such that \(\varrho\) is isomorphic to the representation \(\pi\) of \(\mathbf{R}\) in \(\mathrm{L}^2 (\mathrm{X},\mu)\) defined by \(\pi (t)f = e^{itg}f\) for all \(t\in \mathbf{R}\) and every \(f\in \mathrm{L}^2 (\mathrm{X},\mu)\) .

Let \(u\) the self-adjoint operator on E such that \(\varrho(t) = e^{itu}\) for all \(t \in \mathbf{R}\) (Theorem 1). There exists a locally compact space X, a positive measure \(\mu\) on X, an isometric isomorphism \(\theta\) of \(L^2(X, \mu)\) on E and a continuous function \(g\) on X with real values such that \(u = \theta \circ m_g \circ \theta^{-1}\) (theorem 1 of IV, p. 266). The assertion follows from the formula \(e^{itu} = \theta \circ e^{itm_g} \circ \theta^{-1} = \theta \circ m_{e^{itg}} \circ \theta^{-1}\) (lemma 4 of IV, p. 269).

Remark. --- Suppose that u is an endomorphism of E. The unitary representation of R in E defined by \(\varrho(t) = e^{itu}\) then verifies the inequality \(\|\varrho(t) - 1_{\mathrm{E}}\| \leqslant |t||u\|\) for all \(t \in R\) , and the map \(\varrho\) from R into the Banach space \(\mathcal{L}(\mathrm{E})\) is the unique solution of the linear differential equation

\[
\frac {1}{i} \frac {d \varrho}{d t} = u \circ \varrho
\]

(cf. FVR, IV, p. 26, §6).

Example. --- Let \(\varrho\) the regular representation of \(\mathbf{R}\) in \(\mathrm{L}^2 (\mathbf{R})\) . The infinitesimal generator of \(\varrho\) is the closure of the differential operator with domain \(\mathcal{D}(\mathbf{R})\) defined by \(f\mapsto -if'\) .

\section{FUNCTIONS OF POSITIVE TYPE}

In this section, all vector spaces, as well as all Hilbert spaces and algebras considered, are over C, unless otherwise stated.

\subsection{Universally positive kernels}

In this issue, X is a separated topological space.

Theorem 1. --- Let \(f \in \mathcal{C}(X \times X)\) . The following conditions are equivalent:

\begin{enumerate}
\def\labelenumi{(\roman{enumi})}
\tightlist
\item
  For every compact subset Y of X and every positive measure \(\mu\) on Y, the endomorphism of \(\mathrm{L}^{2}(\mathrm{Y},\mu)\) defined by the kernel \(f|(\mathrm{Y}\times\mathrm{Y})\) (def. 1 of III, p. 29) is positive, in other words, we have
\end{enumerate}

\[
\int_ {\mathrm{Y} \times \mathrm{Y}} \overline {{h (x)}} h (y) f (x, y) d (\mu \otimes \mu) (x, y) \geqslant 0
\]

for all \(h \in \mathcal{L}^{2}(\mathrm{Y}, \mu)\) ;

\begin{enumerate}
\def\labelenumi{(\roman{enumi})}
\setcounter{enumi}{1}
\tightlist
\item
  For every integer \(n \in N\) , every family \((x_{i})_{0 \leqslant i \leqslant n}\) in X and any family \((t_{i})_{0 \leqslant i \leqslant n}\) of complex numbers, we have
\end{enumerate}

\[
\sum_ {i = 0} ^ {n} \sum_ {j = 0} ^ {n} \bar {t} _ {i} t _ {j} f (x _ {i}, x _ {j}) \geqslant 0;
\]

\begin{enumerate}
\def\labelenumi{(\roman{enumi})}
\setcounter{enumi}{2}
\item
  There exists a complex Hilbert space E and a continuous map \(g \colon X \to E\) of total image such that \(f(x, y) = \langle g(x) | g(y) \rangle\) for all \(x\) and \(y\) in \(X\) ;
\item
  There exists a complex Hilbert space E and a continuous map \(g: X \to E\) such that \(f(x, y) = \langle g(x) | g(y) \rangle\) for all x and y in X.
\end{enumerate}

It is clear that (iii) implies (iv), and one sees that (i) implies (ii) by considering the discrete measure \(\mu\) which is the image of the counting measure on \(\{0,\ldots,n\}\) by the mapping \(i\mapsto x_{i}\) and the function h such that

\[
h(x) = \sum_{\substack{0\leqslant i\leqslant n\\ x_{i} = x}}t_{i}.
\]

Let us prove that (iv) implies (i). Suppose there exists a complex Hilbert space E and a continuous map \(g \colon \mathrm{X} \to \mathrm{E}\) such that \(f(x, y) = \langle g(x) \mid g(y) \rangle\) for all \((x, y) \in \mathrm{X} \times \mathrm{X}\) . Let Y be a compact subset of X, \(\mu\) a positive measure on Y and \(h \in \mathcal{L}^{2}(Y, \mu)\) . We have

\[
\begin{array}{r l} & {\int_ {\mathrm{Y} \times \mathrm{Y}} \overline {{h (x)}} h (y) f (x, y) d (\mu \otimes \mu) (x, y)} \\ & {\qquad = \int_ {\mathrm{Y} \times \mathrm{Y}} \overline {{h (x)}} h (y) \langle g (x) | g (y) \rangle d (\mu \otimes \mu) (x, y)} \\ & {\qquad = \left\langle \int_ {\mathrm{Y}} h (x) g (x) d \mu (x) \Big | \int_ {\mathrm{Y}} h (y) g (y) d \mu (y) \right\rangle \geqslant 0} \end{array}
\]

by INT, V, p. 97, § 8, no. 4, prop. 9.

Let us finally prove that (ii) implies (iii). Let \(\widetilde{\mathrm{E}}\) be the space of complex measures with finite support on X. For \(\mu_{1}\) and \(\mu_{2}\) in \(\widetilde{\mathrm{E}}\) we set

\[
\langle \mu_ {1} \mid \mu_ {2} \rangle = \int_ {\mathrm{X} \times \mathrm{X}} f (x, y) (\overline {{\mu}} _ {1} \otimes \mu_ {2}) (x, y).
\]

The sesquilinear form thus defined on \(\widetilde{\mathrm{E}}\) is a positive Hermitian form. Indeed, let \(\mu \in \widetilde{\mathrm{E}}\) ; there exists a finite family \((x_{i})_{0\leqslant i\leqslant n}\) in X and complex numbers \((t_i)_{0\leqslant i\leqslant n}\) such that \(\mu = \sum_{i=0}^{n} t_i \varepsilon_{x_i}\) . We then have

\[
\langle \mu | \mu \rangle = \sum_ {i = 0} ^ {n} \sum_ {j = 0} ^ {n} \bar {t} _ {i} t _ {j} f (x _ {i}, x _ {j}) \geqslant 0
\]

by hypothesis. One defines \(\widetilde{g}\colon\mathbf{X}\to\widetilde{\mathbf{E}}\) by \(\widetilde{g}(x)=\varepsilon_{x}\) . The image of the map \(\widetilde{g}\) generates \(\widetilde{\mathbf{E}}\) . Moreover, for every \((x,y)\in\mathbf{X}\times\mathbf{X}\) , one has on the one hand \(f(x,y)=\langle\widetilde{g}(x)|\widetilde{g}(y)\rangle\) and on the other hand

\[
\| \widetilde {g} (x) - \widetilde {g} (y) \| ^ {2} = f (x, x) + f (y, y) - f (x, y) - f (y, x),
\]

which implies that \(\widetilde{g}\) is continuous since \(f\) is continuous.

Let E be the separated completion Hilbert space of \(\widetilde{E}\) (EVT, V, p. 8, cor. of prop. 4) and \(g\colon X \to E\) the composition of \(\widetilde{g}\) and of the canonical map \(\widetilde{E} \to E\) . Then the map g is continuous, its image is total in E, and one has \(f(x,y) = \langle g(x) | g(y) \rangle\) for all \((x,y) \in X \times X\) .

The method used to prove that (ii) implies (iii) is called the Gelfand-Naimark-Segal construction.

DEFINITION 1. --- One says that a function \(f \in \mathcal{C}(\mathrm{X} \times \mathrm{X})\) is a universally positive kernel on X if it satisfies the equivalent conditions of Theorem 1.

If f is a universally positive kernel on X, a pair \((E,g)\) satisfying condition (iv) of loc. cit. is called a Hilbertian realization of f; if condition (iii) is satisfied, one says that it is a cyclic Hilbertian realization.

We denote \(\mathrm{Noy}_{+}(\mathrm{X})\) the set of universally positive kernels on X.

Let X' be a separated topological space and \(h: X \to X'\) a continuous map. The map \(f \mapsto f \circ (h, h)\) of \(\mathcal{C}(X' \times X')\) in \(\mathcal{C}(X \times X)\) induces by passage to subspaces a map \(\mathrm{Noy}_{+}(X') \to \mathrm{Noy}_{+}(X)\) .

PROPOSITION 1. --- Let f be a universally positive kernel on X. Let \((\mathrm{E}_{1}, g_{1})\) and \((\mathrm{E}_{2}, g_{2})\) be Hilbertian realizations of f. One supposes that \((\mathrm{E}_{1}, g_{1})\) is a cyclic Hilbertian realization. There then exists a unique continuous linear map \(u: E_{1} \to E_{2}\) such that \(g_{2} = u \circ g_{1}\) . This map is isometric. If \((\mathrm{E}_{2}, g_{2})\) is also cyclic, then u is an isomorphism.

The uniqueness of \(u\) follows from the fact that the image of \(g_{1}\) is total in \(\mathrm{E}_{1}\) . Let \(\mathrm{F} = \mathbf{C}^{(\mathrm{X})}\) and let \((e_x)_{x\in \mathrm{X}}\) the canonical basis of \(\mathrm{F}\) . For \(j = 1\) and \(j = 2\) , denote \(u_{j}\) the linear map from \(\mathrm{F}\) in \(\mathrm{E}_{j}\) determined by \(u_{j}(e_{x}) = g_{j}(x)\) , and denote by \(\mathrm{F}_{j}\) its image; the space \(\mathrm{F}_{1}\) is dense in \(\mathrm{E}_{1}\) . Let \(t = \sum t_{x}e_{x}\) be an element of \(\mathrm{F}\) . We have

\[
\begin{array}{c} \| u _ {1} (t) \| ^ {2} = \sum_ {x, y} \bar {t} _ {x} t _ {y} \langle g _ {1} (x)   |   g _ {1} (y) \rangle = \sum_ {x, y} \bar {t} _ {x} t _ {y} f (x, y) \\ = \sum_ {x, y} \bar {t} _ {x} t _ {y} \langle g _ {2} (x)   |   g _ {2} (y) \rangle = \| u _ {2} (t) \| ^ {2}. \end{array}
\]

Consequently, there exists a linear isometric map v from \(F_{1}\) in \(F_{2}\) such that \(u_{2}=v\circ u_{1}\) , and in particular \(g_{2}(x)=v(g_{1}(x))\) for all \(x\in X\) . Since the image of \(g_{1}\) is total in \(E_{1}\) , this map extends to an isometric linear map u from \(E_{1}\) in \(E_{2}\) such that \(g_{2}=u\circ g_{1}\) .

By Lemma 8 of I, p. 107, the image of \(u\) is closed in \(\mathrm{E}_2\) . If \((\mathrm{E}_2,g_2)\) is a cyclic Hilbert realization, the image of \(u\) is also dense in \(\mathrm{E}_2\) and \(u\) is therefore an isomorphism.

PROPOSITION 2. --- The set Noy \(_{+}\) (X) is a self-adjoint cone in the space \(\mathcal{C}(\text{X} \times \text{X})\) ; it is stable under product and it is closed when one equips \(\mathcal{C}(\text{X} \times \text{X})\) with the topology of simple convergence.

It is elementary that if \(f \in \mathrm{Noy}_{+}(\mathrm{X})\) , then \(tf \in \mathrm{Noy}_{+}(\mathrm{X})\) for every real number \(t \geqslant 0\) , and \(\overline{f} \in \mathrm{Noy}_{+}(\mathrm{X})\) .

If \((\mathrm{E}_{1}, g_{1})\) and \((\mathrm{E}_{2}, g_{2})\) are Hilbertian realizations of universally positive kernels \(f_{1}\) and \(f_{2}\) on X, then the pair \((\mathrm{E}_{1} \oplus \mathrm{E}_{2}, g_{1} + g_{2})\) (resp. the pair \((\mathrm{E}_{1} \widehat{\otimes}_{2} \mathrm{E}_{2}, g_{1} \otimes g_{2})\) ) is a Hilbertian realization of \(f_{1} + f_{2}\) (resp. of \(f_{1}f_{2}\) ); they are therefore universally positive kernels.

The characterization (ii) of \(\mathrm{Noy}_{+}(\mathrm{X})\) (theorem 1) implies that this set is closed in \(\mathcal{C}(\mathrm{X}\times\mathrm{X})\) endowed with the topology of simple convergence.

\subsection{Further remarks on the holomorphic functional calculus}

For every subset X of C, we denote by \(X^{*}\) the image of X under complex conjugation. Let U be an open subset of C and \(g: U \to C\) a holomorphic function. The function \(f^{*}: z \mapsto \overline{g(\overline{z})}\) is then defined and holomorphic on \(U^{*}\) . The map \(f \mapsto f^{*}\) is a continuous bijection from \(\mathcal{O}(U)\) to \(\mathcal{O}(U^{*})\) such that \((f_{1}f_{2})^{*} = f_{1}^{*}f_{2}^{*}\) and \((f_{1} + f_{2})^{*} = f_{1}^{*} + f_{2}^{*}\) for \(f_{1}\) and \(f_{2}\) in \(\mathcal{O}(U)\) .

Let C be a compact subset of C. The maps \(f \mapsto f^{*}\) of \(\mathcal{O}(U)\) in \(\mathcal{O}(U^{*})\) as U ranges over the open subsets of C containing C, induce a continuous bijection of the space \(\mathcal{O}(C)\) to \(\mathcal{O}(C^{*})\) (I, p. 49, no. 1), which is also denoted \(f \mapsto f^{*}\) and which satisfies \((f_{1}f_{2})^{*} = f_{1}^{*}f_{2}^{*}\) and \((f_{1} + f_{2})^{*} = f_{1}^{*} + f_{2}^{*}\) for \(f_{1}\) and \(f_{2}\) in \(\mathcal{O}(C)\) .

PROPOSITION 3. --- Let A be a unital involutive Banach algebra. Let \(a \in \mathrm{A}\) . The spectrum of \(a^{*}\) is the image \(\mathrm{Sp}_{\mathrm{A}}(a)^{*}\) of \(\mathrm{Sp}_{\mathrm{A}}(a)\) by complex conjugation. For every \(f \in \mathcal{O}(\mathrm{Sp}_{\mathrm{A}}(a))\) , we have \(f(a)^{*} = f^{*}(a^{*})\) .

The first assertion follows from I, p. 97. By the preceding, the map \(\varphi\) of \(\mathcal{O}(\mathrm{Sp}_{\mathrm{A}}(a))\) into A defined by \(f \mapsto (f^{*}(a^{*}))^{*}\) is a continuous unital morphism of \(\mathcal{O}(\mathrm{Sp}_{\mathrm{A}}(a))\) in A such that the image of the germ in a neighborhood of \(\mathrm{Sp}_{\mathrm{A}}(a)\) of the identity function of \(\mathbf{C}\) is equal to \(a\) . Consequently, \(\varphi\) is the map \(f \mapsto f(a)\) of the holomorphic functional calculus (I, p. 74, th. 5).

Lemma 1. --- Let D be the open unit disk in C. There exists a unique holomorphic function f defined on D such that \(f(z)^{2} = 1 - z\) for all \(z \in D\) and \(f(0) = 1\) . We have \(f^{*} = f\) .

The radius of convergence of the power series

\[
\sum_ {n = 0} ^ {+ \infty} \binom {1 / 2} {n} (- z) ^ {n}
\]

is 1 and its sum \(f\) defines a holomorphic function on D (VAR, R1, p. 27, 3.2.9) taking the value 1 at 0. It satisfies \(f(x)=\sqrt{1-x}\) for all \(x \in D \cap R\) (FVR, III, p. 19), hence \(f(z)^{2} = 1 - z\) for \(z \in D\) since the difference \(f(z)^{2} - (1 - z)\) is a holomorphic function all of whose successive derivatives vanish at 0 (VAR, R1, p. 27, 3.2.5). By definition, one verifies that \(f^{*} = f\) .

Let \(g\) a holomorphic function on D such that \(g(z)^2 = 1 - z\) for all \(z \in \mathrm{D}\) and such that \(g(0) = 1\) . The function \(g\) does not vanish and the continuous function \(f / g\) on D takes values in \(\{-1, 1\}\) ; since D is connected and \(f(0) = g(0)\) , we have \(f = g\) .